\documentclass[11pt,reqno]{amsart}
\usepackage[T1]{fontenc}
\usepackage{lmodern,amsmath,amssymb,mathtools,microtype,booktabs,array,longtable}
\usepackage[margin=1in]{geometry}
\usepackage{graphicx}
\usepackage{xr-hyper}
\usepackage[hidelinks]{hyperref}
\ifdefined\pdfinfoomitdate\pdfinfoomitdate=1\fi
\ifdefined\pdftrailerid\pdftrailerid{}\fi
\ifdefined\pdfsuppressptexinfo\pdfsuppressptexinfo=15\fi
\numberwithin{equation}{section}
\newtheorem{theorem}{Theorem}[section]
\newtheorem{proposition}[theorem]{Proposition}
\newtheorem{lemma}[theorem]{Lemma}
\newtheorem{corollary}[theorem]{Corollary}
\theoremstyle{definition}\newtheorem{definition}[theorem]{Definition}
\newtheorem{example}[theorem]{Example}
\theoremstyle{remark}\newtheorem{remark}[theorem]{Remark}
\newcommand{\E}{\mathbb E}\newcommand{\Prb}{\mathbb P}
\newcommand{\R}{\mathbb R}\newcommand{\one}{\mathbf1}
\newcommand{\Q}{\mathcal Q}\newcommand{\TV}{\mathrm{TV}}
\newcommand{\cp}{\mathrm{cp}}\newcommand{\ev}{\mathrm{ev}}\newcommand{\st}{\mathrm{st}}
\DeclareMathOperator{\tr}{tr}\DeclareMathOperator{\dist}{dist}
\DeclareMathOperator{\law}{Law}\DeclareMathOperator{\rank}{rank}
\DeclareMathOperator{\supp}{supp}
\emergencystretch=2em

\externaldocument[][nocite]{supplement-refs}[General_Theta_Foundations_I_R74_Supplement.pdf]
\hypersetup{pdftitle={General Theta Foundations I: Acquired Geometry and Causal Resource Transfer},pdfauthor={Qian Qi}}
\title[Acquired geometry and causal resource transfer]{General Theta Foundations I:\\Acquired Geometry and Causal Resource Transfer}
\author{Qian Qi}\date{October 2026}
\subjclass[2020]{Primary 60J10, 68Q45; Secondary 20C15, 52A20, 60G40, 62C20, 81P15, 90C22, 94A29}
\keywords{causal experiment, predictive state, positive realization, minimal memory, finite group, quantum instrument, convex duality}
\begin{document}
\begin{abstract}
We study the exact classical memory required to reproduce the laws of a
causal experiment. For compactly reversible finite-dimensional responses,
every finite stochastic realization reduces to a permutation core and at
most one counted initializer. We prove that, for every finite observable
group, sufficiently small positive terminal contrast makes the least core
equal to the least dimension of a permutation module containing the
observable module. All memory thresholds
along the contrast line are attained. For groups whose subgroups are all normal,
an absorption theorem reduces the exact optimization to quotient linear
programs with explicit payoff duals and rational thresholds for rational
data. Sharp examples show positivity transitions invisible to the group
and observable module. The retained theory supplies a common exact classical simulator for general
finite quantum instruments under a uniform
interior condition on every posterior. We strengthen the distinction
between this condition and average mixing: a stationary
qubit family has an exactly depolarizing unconditioned channel,
one-dependent observations and Hankel rank three, but its minimum memory
is exactly $q$ with mixed initialization and $q+1$ with a counted pure start
when its conditional phase has finite order $q\ge3$; a non-root-of-unity
phase requires infinite memory. These results connect
predictive compression, positive realization and finite-resource causal
comparison. Complete stopping, variational, dual, reset, rank-constrained
discrimination, spectral and geometric results retain their separate
hypotheses and proofs.
\end{abstract}
\maketitle
\section{Introduction}\label{sec:current-overview}\label{sec:introduction}
A predictive state records everything about a past experiment that can
affect its future laws. Its linear dimension need not measure the memory
required to realize those laws with nonnegative transition probabilities.
The distinction matters in causal experiments: commands can change what is
observed next, observation changes the state on which later commands act,
and the resources used to retain or communicate information must themselves
be counted. This paper studies that relation between predictive compression
and exact finite realization, within a broader theory of causal resource
comparison and optimal discrimination.

Two forms of observation will be used. A terminal response assigns a
probability row $b(w)$ to every finite command word $w$. An instrument
response assigns a joint probability to every finite command--outcome
history. A classical realization has one fixed collection of stochastic
tables and a finite set of live labels. Unless a mixed-start comparison is
explicitly stated, the initial label is pure and included in the count.
The tables may use arbitrary real probabilities; a lower bound therefore
cannot be evaded by increasing numerical precision. Their description and
execution costs are separate resources.

\subsection*{Reversible responses and the positivity gap}
Let $V_b$ be the span of residual functions $b^w(z)=b(wz)$, with dimension
$d$. Suppose that the command shifts are invertible on $V_b$ and generate
a group with compact closure $G_b$. This compact reversibility condition
holds for prepared finite-dimensional unitary command families after all
directions invisible to future tests have been removed. The finite
realization criterion, Theorem~\ref{thm:r72-reversible}, says that an exact
finite classical realization exists precisely when $G_b$ is finite.

Theorem~\ref{thm:r73-normal-form} gives the resource structure of every
such realization. It replaces continuing stochastic dynamics by a finite
permutation action and at most one transient initializer, without increasing
the number of live labels. The original transition matrices need not form
a group. The proof compresses their observable action through recurrent
classes of words which act as the identity on $V_b$. If $N$ is the number
of deterministic residuals and $\kappa_G(b)$ the minimum size of a
probability factorization $b(g)=\lambda P_gD$, then
\[
 m(b)=\min\{N,1+\kappa_G(b)\}.
\]
The bounded decoder and the initial mixture are essential. The least size
$\nu_G(V_b)$ of a permutation module containing $V_b$ is a lower bound
for $\kappa_G(b)$, and can be strictly smaller.

We determine how this positivity gap changes under terminal readout
contrast. Let $\pi$ be the group-averaged probability row and put
$b_t=\pi+t(b-\pi)$, $0\le t\le1$. For every finite group and every
finite outcome set, Theorem~\ref{thm:r74-low-contrast} proves
\[
 \kappa_G(b_t)=\nu_G(V_b),\qquad
 m(b_t)=\min\{N,1+\nu_G(V_b)\}
\]
for all sufficiently small positive $t$. Throughout $t>0$, the group,
module and residual cardinality stay unchanged. The proof constructs a
strictly positive preparation in a smallest containing permutation module
and controls the decoder explicitly. Theorem~\ref{thm:r74-contrast-thresholds}
then gives a complete finite filtration by attained thresholds. Its general
algebraic computability is distinct from a model-specific closed form.

For every group whose subgroups are normal, including all finite abelian
groups, Theorem~\ref{thm:r74-quotient-lp} makes the optimization linear.
Within-orbit preparation can be absorbed into a quotient decoder; only a
mixture among distinct quotient orbits remains. Theorem~\ref{thm:r74-quotient-dual}
gives necessary and sufficient payoff inequalities, strict exclusion
certificates and exact contrast endpoints. Rational data have rational
optimal controllers and rational thresholds. For abelian groups, character
support determines the small-contrast module bound. At higher contrast,
positivity adds a genuine obstruction: a four-command binary family has
$d=6$ and sixteen residuals at every $t>0$, but its minimum is eleven
labels for $t\le3/5$ and thirteen for $t>3/5$. Explicit payoff witnesses also
give a robust neighbourhood of the larger-memory regime. The retained
prepared-qubit $S_3$ model has an exact four-to-six transition at $2/3$.

\subsection*{Observed instruments and conditional memory}
Reversibility is only one way in which response memory is constrained.
The retained Theorem~\ref{thm:r68-instrument} treats known finite quantum instruments
on $\mathbb C^h$ whose every positive normalized branch state is at least
$\eta I/h$. For one fixed known prepared state it constructs a simulator
with at most
\[
 1+\left\lceil(4h/\eta)^{h-1}\right\rceil
\]
live labels. This theorem was proved before the present revision, together with sparse
branch rows, several-step interiority and a state-positive noise exponent.
A finite physical dictionary contains the whole posterior
region; branchwise barycentric expansions intertwine each completely
positive map with a nonnegative outcome matrix. Consequently the same
tables reproduce all adaptive visible histories, terminal POVMs, almost surely finite stopped
laws and nonnegative history costs. No state-independent outcome
probabilities or commuting-command assumption is required. Applying
depolarizing noise after each observed branch supplies the hypothesis for
arbitrary finite instruments. A complementary common-contraction theorem
gives an explicit finite positive realization of supplied linear responses
under one norm certificate for all commands.

Average mixing is insufficient for the instrument conclusion.
Theorem~\ref{thm:r74-depolarizing-obstruction} constructs a stationary
qubit family with an exactly depolarizing unconditioned channel, uniformly
positive outcome probabilities, one-dependent observations and Hankel rank
three. If its conditional phase has order $q\ge3$, its exact classical
memory is $q$ with a free mixed start and $q+1$ with a counted pure start.
A non-root-of-unity phase requires infinite memory. The upper uses one
invariant polygon for all outcomes; the lower uses an observable cylinder
whose relative peripheral spectrum forces a $q$-cycle. Equality in the
cycle bound is rigid enough to rule out a pure $q$-label initialization.
Thus neither linear rank, unconditioned mixing time nor finite dependence
determines the positive realization resource. The exact finite and infinite
cases lie in one fixed physical dimension and one fixed outcome alphabet.

\subsection*{Classical ingredients and contribution}
Finite residual equivalence, stochastic automata, positive realization and
observable operator models are classical
\cite{RabinScott59,Heller65,CarlylePaz71,Jaeger00,ThonJaeger15,MonrasWinter14}.
Schwarz, Flor and Brown developed the structure of nonnegative and
stochastic matrix groups \cite{Schwarz64,Flor69,Brown64,Schwarz62}.
Symmetric cone factorizations and extended formulations
\cite{GPT13,KPT10}, nonnegative-rank perturbations
\cite{BocciCarliniRapallo11}, and restricted nonnegative factorization
\cite{GillisGlineur12} precede the present work. Their general mechanisms
are credited. The assertions here concern the exact observable normal form
with its initial charge, universal attainment of the permutation-module
bound at small contrast, the normal-quotient linear reduction, and the
resulting exact resource transitions.

Invariant-polytope simulation of quantum instruments and adaptive
composition are explicitly developed by Okay~\cite{Okay26}; pure-state
nets and convex approximation of density matrices are established tools
\cite{AubrunSzarek17}. The earlier instrument results of this paper supply the posterior condition,
quantitative physical dictionary and exact resource accounts, with complete
proofs retained in Section~\ref{sec:interior-instruments}. They are used
here as established premises, not counted again as new results.
Finite-dimensional quantum processes without finite classical realizations
already occur in positive-realization theory \cite{Fanizza24}.
Holroyd and Liggett's finitely dependent colorings provide earlier
one-dependent non-HMM processes \cite{HolroydLiggett16}. We claim neither
first phenomenon. The specific assertion is the simultaneous rank-three,
uniformly positive, exactly depolarizing qubit construction and its sharp
$q$, $q+1$, or infinite classification. Classical finite-group recognition,
real-algebraic minimization and finite-word equivalence also remain credited
\cite{Serre07,BabaiBealsRockmore93,KieferWachter14,Tzeng92}.

\subsection*{Relation to the causal resource theory}
The main article retains the full actual-history comparison theorem and
common paid construction. Its lower witness uses the competitor's actual
pre-call histories. The paired supplement retains performed-law
compactness, variational and dual attainment, robust tied optimality,
normalized stopping rigidity, finite-width transfer, rank-constrained
discrimination, finite-order quantum measurements, two-stage reset control,
memory and communication accounts, spectral geometry and optimal strategy
implementation. It also contains the periodic partition, leakage and
arithmetic laws, all with their original hypotheses and complete proofs.
An independently compiled companion retains absolute calibration and
nonintervening directional sensing.

These layers have distinct quantifiers. General causal comparison applies
to its declared experiments; the group theorem classifies compactly
reversible responses; the posterior theorem supplies a sufficient
condition for general instrument families; the exact examples solve their
specified models; and local geometric results remain local. The theorem
and preservation maps make those dependencies explicit. None of these
results is asserted to complete independent A1, A2 or further Theta
construction objectives.

\section{Prepared kernels, acquired task laws and observers}\label{sec:data}
All hidden-state, action and report spaces are standard Borel. A fixed unknown index $\theta\in\Theta$ specifies an actual preparation and positive normalized causal kernels. Positivity means nonnegative measures, not an everywhere strictly positive density. Dependence on the entire physical past may be included in the hidden state. A legal policy sees past visible reports and its admitted randomization, but not $\theta$ or a future audit. Iteration of these kernels defines its finite-prefix laws; the consistent countable path law is their kernel extension.

Each call has a cost $c_i\ge c_*>0$. The decision to make the next call is predictable. In the geometric route, the declared productive-call flag is also predictable; in the primitive-instrument route, productive reports are specified outcomes and can be random. Commitment is permitted only after a productive report, and occurs exactly once almost surely in each model. At that time a specified future audit $Z\in H$ is scored, where $H\subset\R^k$ is nonempty, compact and convex. The prediction is committed before that audit is observed. A prospective joint report--audit kernel is defined at every productive call, including rejected calls. The stopping coin is independent of the audit conditional on the entire current visible report and the past. Write $L=\operatorname{diam}(H)^2$; the case $L=0$ is trivial whenever a legal rule exists.

The core cost $C_\tau$ includes initial preparation charged to the core, failed calls and all productive calls through commitment. A separately executed final audit must be added to the raw budget. This convention is fixed before taking a risk infimum. Physical time, stored program constants and temporary computation are not identified with call count.

\subsection*{Three resource interfaces}
A checkpoint procedure may use its whole precommit history, but must output one element of a fixed dictionary of at most $M$ values. This counts its output width only; it is not a claim of $\log_2 M$ total working bits.

An autonomous observer has an alphabet $S$, $|S|\le M$, and one fixed parameter-independent program. At an intercall boundary its only persistent observer memory is its label $s\in S$. Its next command distribution is a fixed function of $s$. A physical port name, elapsed time, subprogram index, random seed or former report is not an extra free input. During a call its current command, current report and fresh coins are available. Whatever persists to another call must be encoded in $S$. A fixed decoder $r:S\to H$ reads only the committed label.

In the \emph{event-stopped} interface a post-report kernel emits $(s',e)\in S\times\{0,1\}$. An event $e=1$ ends the experiment and retains $s'$ for the decoder; no subsequent command is executed. This public commitment event is not a stored bit used in continuing computation. The same label may be used live on one path and terminal on another. In the \emph{state-stopped} interface a fixed map $h_*:S\to\{0,1\}$ determines halt status. The initial label is live, and type-one labels are globally absorbing terminal labels with no executed command rows. Entry into a terminal label is allowed only after a productive report. Thus live and terminal types are disjoint in this interface, and only in this interface by definition.

Write $R_i^\Theta(N,M)$, $i\in\{\cp,\ev,\st\}$, for the infimum of worst-case expected squared loss over programs with robust almost-sure commitment and $\sup_\theta\E_\theta C_\tau\le N$. An empty legal class has value $+\infty$. A program chosen once at random is not a free persistent seed. The ideal measurable class may have fixed real constants; effective implementation is stated separately in Section~\ref{sec:composition}.

\subsection*{A physically acquired comparison law}
Execute a declared prepared comparison experiment and obtain its joint report--audit probability $K_0(dx,dz)$ on $X\times H$. Its report marginal is $\nu$. The law is genuinely preparable and its specification has its own resource account. It need not be a member of $\Theta$. In particular no finite set of calibration data is silently promoted to a uniform certificate. Set
\begin{equation}\label{eq:derived}
 f(x)=\int zK_0(dz\mid x),\quad
 w(x)=\int\|z-f(x)\|^2K_0(dz\mid x),\quad
 d_r(x)=w(x)+\|f(x)-r\|^2.
\end{equation}
Disintegration gives Borel versions with $f\in H$; null sets may be completed arbitrarily. The profile is
\begin{equation}\label{eq:profile}
 \Q_m(b)=\min_{\substack{C\subset H,\ 1\le|C|\le m\\
                  0\le h\le b,\ \int h\,d\nu=1}}
        \int [w+\dist(f,C)^2]h\,d\nu,
 \qquad b\ge1,\ m\ge1.
\end{equation}
We use $+\infty$ for an infeasible cap or width. The quantities $f,w$ are task coordinates of acquired predictive information, not a complete substitute for the predictive quotient. Their derivation and the quotient construction in Appendix~\ref{sec:quotient} precede any use of local dimension.

\subsection*{Lower certificates on actual histories}\label{sec:lower-cert}
Fix \emph{one} actual model $\theta_*\in\Theta$, independent of the competing program. At each actual reached productive pre-call history $h$ and legal command $a$, write $P_{h,a}(dy,dz)$ for its conditional full current report--audit law. Let $P^0_{h,a}$ be a measurable comparison row on the same report space, with a Borel mark $x=\pi_{h,a}(y)$ of law $\nu$ and
\begin{equation}\label{eq:full-moments}
 \E^0[Z\mid y]=f(\pi_{h,a}y),\qquad
 \E^0[\|Z\|^2\mid y]=w(\pi_{h,a}y)+\|f(\pi_{h,a}y)\|^2.
\end{equation}
These are conditional identities given the \emph{full} report, not just its mark. Versions and inequalities are required on actual reachable histories, not only almost everywhere under an unrelated comparison path law.

Fix $0<l_-\le1\le u_-<\infty$. For all Borel $q:Y\to[0,1]$ and all finite-image Borel $r:Y\to H$, suppose predictable nonnegative defects $a_-(h,a),e_-(h,a)$ obey
\begin{align}
 \int q\,d(P_{h,a})_Y
 &\le u_-\int q\,d(P^0_{h,a})_Y+a_-(h,a),\label{eq:lower-mass}\\
 \int q(y)\|z-r(y)\|^2P_{h,a}(dy,dz)
 &\ge l_-\int q(y)\|z-r(y)\|^2P^0_{h,a}(dy,dz)-e_-(h,a).
 \label{eq:lower-task}
\end{align}
Independent current randomization can be adjoined to both rows: integrating inequalities for each fixed random value proves the same inequalities with the same defects. All command-conditioned lower inequalities hold before averaging the chosen command.

Let $A_\tau$ count productive calls. At this same actual model require, for every legal competing procedure,
\begin{equation}\label{eq:effort}
 \E_* C_\tau\ge\kappa_-\E_* A_\tau,\qquad\kappa_->0.
\end{equation}
At budget $N$ suppose, uniformly over all robustly feasible competitors,
\begin{equation}\label{eq:occupation}
 \E_*\sum_{j\le\tau:\ j\ {\rm productive}}a_-(H_{j-1},a_j)\le A_N,
 \qquad
 \E_*\sum_{j\le\tau:\ j\ {\rm productive}}e_-(H_{j-1},a_j)\le E_N.
\end{equation}
These are actual occupation certificates, not properties merely of the proposed upper program. Uniform row defects give $A_N\le(N/\kappa_-)\bar a_-$ and $E_N\le(N/\kappa_-)\bar e_-$. Lemma~\ref{lem:drift} supplies a nonuniform alternative.

\subsection*{One common upper construction}\label{sec:upper-cert}
Independently, let $\mathcal G$ be a parameter-independent program which, from each of its start histories, reaches one productive report. Its conditional expected attempt cost is at most $\kappa_+>0$ on every surviving history in every $\theta\in\Theta$. It uses $B$ live roles and returns to its start label on rejection; no reset of the physical state is assumed. It forms a virtual mark $x$, a current choice $J$ and a permitted export $T_Jr$ of its task readout. For every finite dictionary $C$ assume
\begin{equation}\label{eq:exports}
 |\{T_jr:j\in J,\ r\in C\}|\le D_0|C|,
 \qquad D_0\in\mathbb N.
\end{equation}
Every continuing phase and coin needed for this construction is included in the $B$ live roles. For fixed $0<l_+\le1\le u_+<\infty$ and $a_+,e_+\ge0$, assume for all Borel $q:X\to[0,1]$ and finite-image Borel $r:X\to H$, conditionally on every such history,
\begin{align}
 \E[q(x)]&\ge l_+\int q\,d\nu-a_+,\label{eq:upper-mass}\\
 \E[q(x)\|Z-T_Jr(x)\|^2]
 &\le u_+\int q(x)d_{r(x)}(x)\nu(dx)+e_+.
 \label{eq:upper-task}
\end{align}
The expectation here includes the within-attempt choice $J$. Unlike the lower, the upper is not required to hold for each choice separately. The choices retained after this call are nevertheless charged through the actual export or live role. These hypotheses express one common causal construction, not modelwise optimizers.

\section{Reversible response groups and common controlled memory}\label{sec:r72-reversible}
The fixed-phase results concern one transition matrix. Several commands require one family of matrices on the same retained labels, and checking every command in isolation is insufficient. We give an intrinsic finite-realization criterion for responses whose minimal linear dynamics is reversible and compact, including any finite family of unitary commands on a prepared finite-dimensional quantum system. We then obtain a matching arithmetic memory law for a pair of noncommuting involutions. This is a theorem about exact controlled response realization, not an assertion that a simulator's private label is an observed sufficient statistic for arbitrary feedback optimization.

Positive realization through invariant polyhedral sets is classical \cite{Heller65,MonrasWinter14,ThonJaeger15}. The argument below uses a section of the realizing simplex before passing to the observable quotient; auxiliary hidden-state responses need not belong to the minimal response space. Compact reversibility then turns containment into equality. The resulting finite-group obstruction concerns all command words, rather than a single spectral restriction.

\subsection{The minimal response space}
Let $A$ and $C$ be nonempty finite alphabets of commands and terminal outcomes. A supplied response is a function $b:A^*\to\Delta(C)$. Its residual after a word $w$ is $b^w(z)=b(wz)$, including the empty word $\varnothing$. Put
\[
 V_b=\operatorname{span}_{\R}\{b^w:w\in A^*\},
 \qquad (L_af)(z)=f(az).
\]
Assume $V_b$ is finite dimensional. It is invariant under each $L_a$, since $L_ab^w=b^{wa}$. Call the response \emph{compactly reversible} if the restrictions $L_a|_{V_b}$ are invertible and the closure $G_b$ of their generated group is a compact subgroup of $\operatorname{GL}(V_b)$. This concerns the observable quotient; it does not require the full physical command group to be finite or fully observable.

An exact stochastic realization with $M$ live labels has row-stochastic matrices $Q_a$, one row-stochastic terminal decoder matrix $D$, and the prescribed initial row $e_1^{\mathsf T}$, such that
\begin{equation}\label{eq:r72-response}
 b(a_1\cdots a_n)=e_1^{\mathsf T}Q_{a_1}\cdots Q_{a_n}D.
\end{equation}
Randomized terminal readout is allowed even when all $Q_a$ are deterministic. The target input word is not encoded into an uncounted clock or persistent seed.

\begin{theorem}[Finite realization of compactly reversible responses]\label{thm:r72-reversible}
For a compactly reversible response, the following are equivalent:

\textup{(i)} An exact finite stochastic realization exists.

\textup{(ii)} The residual set $\{b^w:w\in A^*\}$ is finite.

\textup{(iii)} The observable group $G_b$ is finite.

\textup{(iv)} An exact realization with deterministic transitions and a possibly randomized decoder exists.

The least number of deterministic-transition labels is the residual cardinality. This is not in general the least stochastic label count. If an $M$-label stochastic realization exists, then
\[
 |G_b|\le(2^M-1)!.
\]
This last bound is only a finite combinatorial certificate, not a sharp memory law.
\end{theorem}
\begin{proof}
Suppose \eqref{eq:r72-response} holds, and write $r_w=e_1^{\mathsf T}Q_w$. Let $R=\operatorname{span}_{\R}\{r_w\}\subseteq\R^M$. The assignment $\pi(r_w)=b^w$ extends to a well-defined surjective linear map $\pi:R\to V_b$: a relation $\sum_wc_wr_w=0$ implies, after right multiplication by $Q_zD$, the identical relation among $b^w(z)$ for every $z$. Moreover
\[
 \pi(rQ_a)=L_a\pi(r)\qquad(r\in R).
\]
The simplex section $K=R\cap\Delta_M$ is a compact polytope, invariant under right multiplication by every $Q_a$. Its image $P=\pi(K)$ is a compact polytope, contains every residual, and spans $V_b$. It satisfies $L_aP\subseteq P$. It is important that we intersect the simplex with $R$ first: projecting arbitrary hidden-state responses directly into $V_b$ would not be justified.

For an element $T$ of a compact matrix group, positive powers have a subsequence tending to the identity with exponents tending to infinity. One elementary proof takes a convergent subsequence of powers with successively growing index gaps; quotients of two close powers approach the identity. A finite-order element is covered by its period multiples. Thus positive powers also approach $T^{-1}$. Applying this to $L_a$ and using closedness of $P$ gives $L_a^{-1}P\subseteq P$, and hence $L_aP=P$. All elements of $G_b$ preserve $P$. They permute its finitely many vertices, and this permutation action is faithful because those vertices span $V_b$. Consequently $G_b$ is finite. This proves (i)$\Rightarrow$(iii).

For the displayed bound, a vertex of $K$ is determined by its positive support. If two distinct vertices had the same support, their difference could be added and subtracted in sufficiently small multiples while preserving $R$, unit mass and nonnegativity, contradicting extremality. Thus $K$ has at most $2^M-1$ vertices. A linear image of a polytope has no more vertices than the original polytope, so the faithful vertex action proves the bound.

If $G_b$ is finite, the residual orbit is finite. Conversely, if the residual set is finite, each invertible $L_a$ maps it injectively into itself, hence permutes it. Since it spans $V_b$, the generated group and its closure are finite. This proves (ii)$\Leftrightarrow$(iii). Under (ii), use the residuals as labels, send $b^w$ to $b^{wa}$ on input $a$, and use decoder $b^w(\varnothing)$. This realizes $b$ with deterministic transitions. If any such deterministic-transition machine sends $u$ and $v$ to the same label, every suffix yields the same decoder law, so $b^u=b^v$. This proves its minimal count. Finally (iv)$\Rightarrow$(i) is immediate. No assertion that stochastic transitions can always be determinized without increasing size was used.
\end{proof}

\begin{corollary}[Prepared unitary commands]\label{cor:r72-unitary}
Let $U_a$ be finitely many unitary matrices on $\mathbb C^d$, let $\rho$ be a prepared density matrix, and let $\{E_r:r\in C\}$ be a terminal positive-operator-valued measurement. For
\[
 b(w,r)=\tr(E_r\Phi_w(\rho)),\qquad \Phi_a(X)=U_aXU_a^*,
\]
exact finite stochastic realization for all command words is possible if and only if its observable response group is finite. The group is computed after quotienting only those directions invisible to every future commanded terminal test. Its dimension is at most $d^2$.
\end{corollary}
\begin{proof}
Let $W=\operatorname{span}_{\R}\{\Phi_w(\rho)\}$ in the Hermitian matrices. Each conjugation is a Hilbert--Schmidt isometry. Its forward invariant subspace $W$ is also backward invariant, by finite-dimensional invertibility. Let
\[
 N=\{X\in W:\tr(E_r\Phi_zX)=0\text{ for all }z,r\}.
\]
It is invariant under every $\Phi_a$ and hence under its inverse. The orthogonal complement of $N$ in $W$ is invariant as well. Therefore the induced transformations on $W/N$ are orthogonal for the quotient inner product. The map $X\mapsto(z,r\mapsto\tr(E_r\Phi_zX))$ identifies $W/N$ with $V_b$ and intertwines them with $L_a$. Their generated group has compact closure in the orthogonal group. Apply Theorem~\ref{thm:r72-reversible}. A constant measurement can make this induced group trivial even when the physical unitary group is infinite, so finiteness of the full physical group is not a necessary condition.
\end{proof}

\subsection{A matching memory law for two noncommuting commands}
Fix $x\in\R$ with $\sin(2\pi x)\ne0$. On a protected qubit take
\[
 U_0=\begin{pmatrix}1&0\\0&-1\end{pmatrix},\qquad
 U_1=\begin{pmatrix}\cos(\pi x)&\sin(\pi x)\\
                     \sin(\pi x)&-\cos(\pi x)\end{pmatrix},
 \quad \rho=E=|0\rangle\langle0|.
\]
Both commands are involutions. For known $0\le\alpha\le1$, their channels are
\[
 \Phi_{a,\alpha}(X)=\alpha U_aXU_a^*+(1-\alpha)\tr(X)I/2.
\]
Let $M_x^{\rm ctrl}(\alpha)$ be the least stochastic live memory of the whole binary response $w\mapsto\tr(E\Phi_{w,\alpha}\rho)$, using the same matrices and decoder for every word. In the real Bloch plane, with coordinates $(z,x)$, these commands are $\alpha S_0$ and $\alpha S_1$, where $S_0$ reflects the angle $t$ to $-t$ and $S_1$ reflects it to $2\pi x-t$. Their product is $\alpha^2R_{2\pi x}$.

\begin{theorem}[All-word arithmetic law for an involutive command pair]\label{thm:r72-control-envelope}
For $0<\epsilon<1$ and $\alpha=1-\epsilon$, the arithmetic envelope of Section~\ref{sec:r71-arithmetic} gives
\begin{equation}\label{eq:r72-control-envelope}
 \tfrac14\mathfrak A_x(\epsilon)
 \le M_x^{\rm ctrl}(1-\epsilon)
 \le(1+4\pi^2)\mathfrak A_x(\epsilon).
\end{equation}
The upper is one explicit finite stochastic machine valid for every mixed command word, not separately chosen unary machines. The lower ranges over arbitrary common stochastic realizations.

For irrational $x$, the lower and upper logarithmic growth exponents as $\epsilon\downarrow0$ are respectively
\[
 \frac1{\mu(x)+1},\qquad\frac{\mu(x)-1}{2\mu(x)-1},
\]
with values $0,1/2$ at $\mu=\infty$. A uniform cubic-root law holds precisely for badly approximable $x$; otherwise the normalized count has lower limit zero and upper limit infinity.

At $\alpha=1$, the exact count is the reduced denominator $q$ if $x$ is rational, and is infinite otherwise. At $\alpha=0$, it is two. Every strictly damped pair has a finite exact common realization. These statements concern controlled response memory, not a general multi-command boundary-optimal policy theorem.
\end{theorem}
\begin{proof}
Any common realization restricted to the words $(01)^n$ realizes
\[
 \tfrac12[1+\alpha^{2n}\cos(2\pi nx)]
\]
with transition matrix $Q_0Q_1$. Its size is therefore at least
$m_x(1-\alpha^2)\ge\mathfrak A_x(\epsilon(2-\epsilon))/2$
by Theorem~\ref{thm:r71-envelope}. For $1\le c\le2$, the defining maximum of the envelope obeys $\mathfrak A_x(c\epsilon)\ge\mathfrak A_x(\epsilon)/c$ whenever both arguments are in its domain. This follows for each denominator before taking its infimum. Taking $c=2-\epsilon$ proves the lower in \eqref{eq:r72-control-envelope}.

For the upper, the rational multiple-polygon construction in the proof of Theorem~\ref{thm:r71-envelope} gives a regular $l$-gon $P$ on the unit circle, with vertex $v_0=(1,0)$, such that
\[
 l\le(1+4\pi^2)\mathfrak A_x(\epsilon),\qquad
 \alpha R_{2\pi x}v_j\in P\quad\text{for every vertex }v_j.
\]
The first reflection $S_0$ permutes this regular polygon and zero belongs to it. Consequently $\alpha S_0P\subseteq P$, while $S_1=R_{2\pi x}S_0$ gives $\alpha S_1P\subseteq P$. Choose barycentric rows $Q_a(j,k)$ of $\alpha S_av_j$ in the same vertices. Let $d_j=(1+(v_j)_1)/2$. Starting at vertex zero, induction on the actual command word shows that the expected retained vertex equals the protected Bloch vector. Pairing with the first coordinate proves every response identity. The start is one of the counted vertices; there is no extra initialization mixture.

The two constants in \eqref{eq:r72-control-envelope} transfer both exponents and both uniform-law statements from Theorems~\ref{thm:r71-exponents} and \ref{thm:r71-cubic-rigidity}. At zero damping, a rational phase of reduced denominator $q\ge3$ has a common deterministic $q$-gon realization. The word-pair transition has a primitive $q$th-root mode and needs at least $q$ states by the unit-circle spectral obstruction. For irrational $x$ that same obstruction rules out every finite size. At $\alpha=0$ the response is one at the empty word and $1/2$ after every nonempty word, realized by a first label followed by an absorbing fair label; one label cannot change its decoder. This proves all endpoints.
\end{proof}

A fully algebraic example uses $U_1=5^{-1/2}\left(\begin{smallmatrix}2&1\\1&-2\end{smallmatrix}\right)$. Then $e^{2\pi ix}=(3+4i)/5$ is not a root of unity, since its sum with its conjugate is the rational noninteger $6/5$. Each undamped command in isolation has period at most two; their common all-word response has no finite stochastic realization. The arithmetic exponent of this particular angle is not determined here. The example isolates the failure of commandwise testing without inferring a Diophantine class from finitely many coefficients.

If an independent herald with fixed probability $p$ is placed before each command operation, and only failures apply the supplied command, every common realization above preserves the terminal response under every external command schedule and the first-herald cost $1/p$. The commands are actual inputs to that same retained-state transducer. This does not make its private label observable to the caller, allow uncharged midstream quantum measurements, or prove optimality for a controller free to choose an easier command sequence. The unary optimal squared-task interpretation below continues to use its separately proved physical embedding.

\section{A permutation normal form with an exact memory charge}
\label{sec:r73-normal-form}

Theorem~\ref{thm:r72-reversible} identifies when finite memory exists. We now
determine the form of every possible finite realization, retaining the cost of
its prescribed pure initial label. The distinction between a minimal linear
representation and a positive realization remains essential. In particular,
the permutation representation carried by the observable space does not by
itself provide a decoder with values in the probability simplex.

Let $b$ be a compactly reversible response admitting a finite realization.
Write $G$ for its finite observable group, $d=\dim V_b$, and $N$ for its residual
cardinality. We use the chronological convention $g(u)g(v)=g(uv)$: the right
action on $V_b$ is $f\cdot g(w)=L_wf$, where $L_wf(z)=f(wz)$. Thus $G$ in this
notation is the opposite of the group of operators with the usual convention
for composition; inversion identifies the two groups. No group cardinality,
faithfulness, or representation-theoretic statement depends on this convention.
Every element has a representative in $A^*$, including its inverse, since a
finite group generated by the command elements is already generated by their
positive words. The value $b(g)=b(w)$ is consequently well defined.

A \emph{probability factorization on a finite $G$-set} consists of a finite set
$X$, row permutation matrices $P_g$ with $P_gP_h=P_{gh}$, a probability row
$\lambda$ on $X$, and a row-stochastic matrix $D$ with columns indexed by $C$,
such that
\begin{equation}\label{eq:r73-probability-factorization}
 b(g)=\lambda P_gD\qquad(g\in G).
\end{equation}
All rows of $D$ are probability laws, including rows associated with states
that are not distinguished by the immediate terminal measurement. Unused
orbits may be deleted. Define
\[
 \kappa_G(b)=\min\{|X|:\text{\eqref{eq:r73-probability-factorization} holds}\}.
\]
The residual realization shows that this minimum is finite and at most $N$.
It permits a mixed $\lambda$ without charging for its preparation; the
stochastic memory $m(b)$ continues to charge every live label, beginning at
the specified point mass as in \eqref{eq:r72-response}.

\begin{theorem}[Reversible stochastic normal form]\label{thm:r73-normal-form}
Every $M$-label stochastic realization of $b$ gives a faithful finite $G$-set
$X$ and a factorization \eqref{eq:r73-probability-factorization} satisfying
\begin{equation}\label{eq:r73-initialization-charge}
 |X|+\mathbf 1_{\{\lambda\text{ is not a point mass}\}}\le M.
\end{equation}
Conversely, every such factorization has a realization with the displayed
number of labels. Its transitions on $X$ are permutations; if $\lambda$ is
mixed, its only additional label initializes the distribution on the first
command. Thus arbitrary continuing stochastic dynamics confers no advantage
over one initialization mixture followed by a permutation action.
\end{theorem}

\begin{proof}
Fix a realization $(e_1^{\mathsf T},\{Q_a\},D_0)$, and retain the reachable
row space $R$ and the surjection $\pi:R\to V_b$ from the proof of
Theorem~\ref{thm:r72-reversible}. Write
\[
 \mathcal L=\{w\in A^*:g(w)=1\}.
\]
These are identity words in the observable group; their matrices need not be
identity matrices on hidden states. Choose finitely many words
$\ell_1,\ldots,\ell_s\in\mathcal L$ whose positive matrix entries cover all
directed edges occurring in any $Q_\ell$, $\ell\in\mathcal L$. There are at
most $M^2$ distinct edges, so $s\le M^2$ is possible. If no nontrivial choice
is needed, the empty word suffices. Set
\[
 T=\frac{I+Q_{\ell_1}+\cdots+Q_{\ell_s}}{s+1}.
\]
Its closed communicating classes $C_1,\ldots,C_k$ are closed for every
$Q_\ell$ with $\ell\in\mathcal L$. The positive diagonal makes the
restriction of $T$ to each such class aperiodic. The elementary convergence
argument in Lemma~\ref{lem:r73-lazy-convergence} gives a stochastic idempotent
\[
 E=\lim_{n\to\infty}T^n,
 \qquad
 \Delta_ME=\operatorname{conv}\{\mu_1,\ldots,\mu_k\},
\]
where $\mu_i$ is the unique stationary probability row of $T$ on $C_i$.
Their supports are disjoint, so this image is a simplex with exactly $k$
vertices. Every identity-word matrix preserves each $C_i$; hence
\begin{equation}\label{eq:r73-sandwich}
 EQ_\ell E=E\qquad(\ell\in\mathcal L).
\end{equation}
This is a sandwiched equality. We do not assert $Q_\ell E=E$, $EQ_\ell=E$,
or common stationarity of the $\mu_i$ under all identity words.

Each $T^n$ is a convex combination of identity-word matrices. Consequently
$RE\subseteq R$ and $\pi(rE)=\pi(r)$ for $r\in R$. For a word $u$, put
$F_u=EQ_uE$, considered as a stochastic map of $\Delta_ME$ to itself.
Choose a word $\bar u$ representing $g(u)^{-1}$. Replacing the middle $E$
by the convex combinations $T^n$ and using
$u\ell\bar u\in\mathcal L$ gives
\[
 F_uF_{\bar u}=E=F_{\bar u}F_u.
\]
Thus $F_u$ is an affine automorphism of that simplex. Such an automorphism
permutes its vertices, because both it and its inverse preserve the simplex
and therefore preserve extremality. If $g(u)=g(v)$, the same inverse word
shows $F_uF_{\bar u}=F_vF_{\bar u}=E$, so $F_u=F_v$. Next,
$EQ_uT^nQ_vE$ is a convex combination of $F_{u\ell v}$, all equal to
$F_{uv}$. Passing to the limit proves $F_uF_v=F_{uv}$. We have obtained a
permutation action $P_g$ on the vertices $X=\{\mu_1,\ldots,\mu_k\}$.

This action is faithful. If $F_u$ is the identity on $\Delta_ME$, then
$rEQ_uE=rE$ for all rows $r$. For $r\in R$, application of $\pi$ shows
$L_u\pi(r)=\pi(r)$. Surjectivity gives $L_u=I$ on $V_b$, and so $g(u)=1$.
Write $e_1^{\mathsf T}E=\sum_i\lambda_i\mu_i$ and set
$D(i,\cdot)=\mu_iD_0$. Intertwining on $R$ gives, for every word $w$,
\[
 \lambda P_{g(w)}D=e_1^{\mathsf T}EQ_wED_0=b(w),
\]
including the empty word. If the original initial label lies in $C_i$, its
row under $E$ is exactly $\mu_i$. A mixed $\lambda$ therefore requires that
the initial label lie outside the union of the closed classes. The $k$
nonempty disjoint classes then use at most $M-1$ labels. This proves
\eqref{eq:r73-initialization-charge}; a pure $\lambda$ only requires $k\le M$.

For the converse, when $\lambda$ is a point mass its support point is the
initial label in $X$. If it is mixed, adjoin a label $\star$ with decoder
$D(\star,\cdot)=\lambda D$. On command $a$, its transition row into $X$ is
$\lambda P_{g(a)}$, and no transition returns to $\star$. The rows on $X$
are those of $P_{g(a)}$. Induction gives \eqref{eq:r73-probability-factorization}
for every nonempty word, and the decoder at $\star$ gives the empty word.
The first input is applied during initialization, so it is not lost or
encoded in an uncounted preparatory clock.
\end{proof}

The general structure of nonnegative idempotents and bounded groups of
nonnegative matrices is classical; see Brown~\cite{Brown64},
Schwarz~\cite{Schwarz64}, and Flor~\cite{Flor69}. The proof above supplies the
additional steps required for an arbitrary semigroup realization: the
identity-word projection is constructed before passing to the observable
quotient, the sandwiched action is proved to factor through that quotient,
and the prescribed pure start is charged. It does not assume that the
original stochastic matrices themselves form a group.

\begin{theorem}[Exact memory formula]\label{thm:r73-memory-formula}
With randomized terminal readout allowed and the initial label prescribed
pure, the minimum stochastic memory is
\begin{equation}\label{eq:r73-memory-formula}
 m(b)=\min_{(X,\lambda,D)}
 \left(|X|+\mathbf1_{\{\lambda\text{ is not a point mass}\}}\right)
 =\min\{N,1+\kappa_G(b)\}.
\end{equation}
The minimum ranges over precisely the probability factorizations
\eqref{eq:r73-probability-factorization}, not arbitrary linear factorizations.
\end{theorem}
\begin{proof}
The first equality is both directions of Theorem~\ref{thm:r73-normal-form}.
A point-mass factorization gives a deterministic-transition realization on
$X$, so $|X|\ge N$ by Theorem~\ref{thm:r72-reversible}. A mixed factorization
costs at least $1+\kappa_G(b)$. Conversely, the residual machine has $N$
labels, and a factorization attaining $\kappa_G(b)$ costs at most
$1+\kappa_G(b)$. These inequalities give the second equality, also when a
minimum factorization happens to have a pure start.
\end{proof}

\begin{corollary}[Sharp group bounds]\label{cor:r73-group-bound}
An $M$-label realization of a compactly reversible response satisfies
\[
 |G|\le M!,\qquad
 M<N\ \Longrightarrow\ |G|\le(M-1)!.
\]
Both bounds are attained by families below. In particular, a stochastic
memory saving over the residual count requires a faithful permutation core
on at most $M-1$ points.
\end{corollary}
\begin{proof}
The core action in Theorem~\ref{thm:r73-normal-form} embeds $G$ into $S_k$.
If $M<N$, its start cannot be pure, since that would produce a deterministic
machine with at most $M$ labels. Thus $k\le M-1$ in this case.
\end{proof}

Let $\nu_G(V_b)$ be the least cardinality of a finite $G$-set whose real
permutation module contains $V_b$, and let $\mu(G)$ be the least cardinality
of a faithful finite $G$-set. All references to $V_b$ here use the right-action
convention already fixed. A probability factorization gives $V_b$ as a
subquotient of $\mathbb R^X$: take the cyclic span of $\lambda$, then its
observable quotient. Averaging an inner product over finite $G$ splits every
invariant subspace orthogonally, so this subquotient embeds in
$\mathbb R^X$. We obtain
\begin{equation}\label{eq:r73-module-bounds}
 d\le\nu_G(V_b)\le\kappa_G(b),\qquad
 m(b)\ge\max\{d,\mu(G),\min(N,1+\nu_G(V_b))\}.
\end{equation}
These are lower bounds, not an identification of a linear module with a
bounded decoder. Proposition~\ref{prop:r73-decoder-obstruction} makes the
distinction strict.

\begin{corollary}[Rank rigidity]\label{cor:r73-rank-rigidity}
For a finite compactly reversible response,
\[
 m(b)=d\quad\Longleftrightarrow\quad N=d.
\]
In that case every realization with $d$ labels has permutation transition
matrices. If $N>d$, then $m(b)\ge d+1$.
\end{corollary}
\begin{proof}
Every $M$-label realization has $d\le M$. If $M=d$, the core in
Theorem~\ref{thm:r73-normal-form} has $k\ge d$ by its linear rank, hence
$k=d$ and its start is pure. Thus $N\le d$, whereas $N\ge d$ because the
residuals span $V_b$. Conversely $N=d$ gives a deterministic machine with
$d$ labels. For the assertion about the original matrices, a $d$-label
realization has $R=\mathbb R^d$ and the observable map $\pi$ is an
isomorphism. An inverse word for $g(a)$ therefore gives a stochastic inverse
of $Q_a$ on all of $\mathbb R^d$. A stochastic map with a stochastic inverse
is an affine automorphism of the full simplex and hence a permutation.
\end{proof}

\subsection{Two commands and a factorial deterministic gap}
Let $k\ge3$, let $P_g$ be the natural row permutation action of $S_k$, choose
a nonconstant $\lambda\in\Delta_k$, and choose a nonconstant column
$t\in[0,1]^k$. Two commands are the $k$-cycle $(1\,2\,\cdots\,k)$ and the
transposition $(1\,2)$; these generate $S_k$. Supply the binary response
\begin{equation}\label{eq:r73-symmetric-response}
 b(g)=\bigl(\lambda P_gt,\,1-\lambda P_gt\bigr).
\end{equation}
If the distinct coordinate values of $\lambda$ have multiplicities
$n_1,\ldots,n_j$, put $N_\lambda=k!/(n_1!\cdots n_j!)$.

\begin{proposition}[Exact symmetric-group family]\label{prop:r73-symmetric-family}
The response \eqref{eq:r73-symmetric-response} has observable group $S_k$,
dimension $k$, residual cardinality $N_\lambda$, and
\[
 \kappa_{S_k}(b)=k,\qquad
 m(b)=\min\{N_\lambda,k+1\}.
\]
When the coordinates of $\lambda$ are pairwise distinct, the least stochastic
memory is $k+1$ and the least deterministic-transition memory is $k!$.
For a point-mass $\lambda$, both are $k$. If either $\lambda$ or $t$ is
constant, the response instead has one residual and one-label memory.
\end{proposition}
\begin{proof}
For any nonconstant vector, differences between it and its transposes yield
a nonzero multiple of some $e_i-e_j$. Conjugating by all permutations shows
that the span of its orbit differences is
$\mathbf1^\perp$. The orbit of $\lambda$ therefore spans $\mathbb R^k$,
because every row has mass one. The orbit of $t$ also spans
$\mathbb R^k$: its differences span $\mathbf1^\perp$, and its coordinate
sum is positive, since a nonconstant nonnegative vector is nonzero.
Consequently the map
\[
 r\longmapsto\bigl(h\longmapsto(rP_ht,\,r\mathbf1-rP_ht)\bigr)
\]
is injective on $\mathbb R^k$. It identifies the reachable row space with
$V_b$, and identifies the observable action with the faithful natural
permutation action. Its residuals are precisely the distinct permutations
of $\lambda$, giving $d=k$ and $N=N_\lambda$. The displayed factorization
and the rank lower bound give $\kappa=k$. The memory formula proves the
remaining assertion; it also shows $N_\lambda\ge k$. The two constant
cases are immediate from \eqref{eq:r73-symmetric-response}.
\end{proof}

This is an exact prepared quantum experiment in every dimension $k$: take
$\rho=\operatorname{diag}(\lambda)$, permutation unitaries for the commands,
and terminal effect $E=\operatorname{diag}(t)$. The construction already
lies in the diagonal, classical subalgebra. It demonstrates a stochastic
versus deterministic memory separation, and does not assert a quantum
advantage for this family. It also proves sharpness of
Corollary~\ref{cor:r73-group-bound}: the point-mass family has $M=k$ and
$|G|=k!$, while pairwise distinct coordinates have $M=k+1<N=k!$ for $k\ge3$.

For a concrete rational example, take $k=3$,
$\lambda=(1/2,1/3,1/6)$, and $t=(0,1/4,1)^{\mathsf T}$.
The six first-outcome values are
\[
 \left\{\frac14,\frac38,\frac7{24},\frac{13}{24},
                 \frac{11}{24},\frac7{12}\right\}.
\]
It has $d=3$, $N=6$, $\kappa=3$, and $m=4$.

\subsection{The decoder constraint can remove the entire saving}
The representation-theoretic lower bound in
\eqref{eq:r73-module-bounds} need not attain the positive factorization
minimum, even for the same group, rank, and residual count as this example.

\begin{proposition}[A strict bounded-decoder obstruction]
\label{prop:r73-decoder-obstruction}
Let $u=(1,-1,0)$ and let $S_3$ act by its natural permutation matrices. The
binary response with first coordinate
\begin{equation}\label{eq:r73-endpoint-response}
 \beta(g)=\frac12+\frac14 uP_gu^{\mathsf T}
\end{equation}
has $d=3$, $N=6$, and $\nu_{S_3}(V_b)=3$, but
\[
 \kappa_{S_3}(b)=m(b)=6.
\]
It is realized by a prepared qubit with
\[
 U_0=\begin{pmatrix}1&0\\0&-1\end{pmatrix},\qquad
 U_1=\frac12\begin{pmatrix}1&\sqrt3\\\sqrt3&-1\end{pmatrix},
 \qquad \rho=E=|+\rangle\langle+|.
\]
Thus the least stochastic memory can differ for responses having the same
observable group, module, and residual cardinality.
\end{proposition}
\begin{proof}
The orbit of $u$ consists of its six distinct permutations and spans the
two-dimensional zero-sum subspace. Pairing with all permutations of $u$
is nondegenerate on that subspace. Including the constant outcome sum gives
dimension three and six distinct residuals. This representation is the
natural permutation module $\mathbb R^3$, and the rank bound gives
$\nu=3$.

Write $s=(1\,2)$ and $r=(1\,2\,3)$. Directly from
\eqref{eq:r73-endpoint-response},
\[
 \beta(1)=1,\quad\beta(s)=0,\quad
 \beta(r)=\beta(r^2)=\frac14,
 \quad\beta(\tau)=\frac34
 \quad\text{for the other two transpositions }\tau.
\]
Suppose a probability factorization exists on at most five points. Its
action is faithful, by the same residual argument as in
Theorem~\ref{thm:r73-normal-form}. Transitive $S_3$-sets have sizes
$1,2,3,6$, since their stabilizers have orders $6,3,2,1$; the subgroups of
order two are conjugate and the subgroup of order three is unique. Thus
our faithful action has one natural three-point orbit and, at most, one
two-point sign orbit together with fixed points. Two three-point orbits
already exceed the budget, and actions on only fixed and sign orbits
annihilate the alternating subgroup.

Let $a$ be any point of positive initial mass and write $q$ for the
first-coordinate decoder. The endpoint values force $q(a)=1$ and
$q(a\cdot s)=0$. No fixed point can carry mass. Within the natural orbit,
the point fixed by $s$ cannot carry mass, and both of the points exchanged
by $s$ cannot carry positive mass simultaneously. Thus the initial mass in
that orbit is concentrated at one moved point $a$; denote its total by
$A$. If a sign orbit is present its mass is concentrated at a single point;
write this total as $B$. We have $A+B=1$. Every even permutation has decoder
one on the mass in the sign orbit, and every odd permutation has decoder
zero there.

The images $a\cdot r$ and $a\cdot r^2$ are the other two points in the
natural orbit. One is $a\cdot s$, with decoder zero, and the other's decoder
is some $t\in[0,1]$. The two equal values $\beta(r)=\beta(r^2)=1/4$ give
$B=1/4$ and $At=0$. Hence $A=3/4$ and $t=0$. Of the two transpositions
different from $s$, one fixes $a$ and the other sends it to the third point.
Their first-outcome values would therefore be $A$ and zero, whereas both
must be $3/4$. This contradiction rules out all factorizations on at most
five points. The regular six-point residual machine gives $\kappa\le6$,
so $\kappa=6$ and Theorem~\ref{thm:r73-memory-formula} gives $m=6$.

For the physical realization, the adjoint actions of $U_0$ and $U_1$ on
the real Bloch plane are two reflections whose product is a rotation by
$2\pi/3$. The initial and terminal Bloch vector is the $x$-axis vector.
The first reflection sends it to its negative, the other two reflections
have inner product $1/2$ with it, and the two nonidentity rotations have
inner product $-1/2$. The Born probability is one half of one plus that
inner product. It has exactly the displayed six group values, so it is
\eqref{eq:r73-endpoint-response} after the specified generator assignment.
\end{proof}

The cyclic partition formula of the preceding revision is recovered as a
special case in Proposition~\ref{prop:r73-cyclic-module}. For the general
finite group, however, both the start charge and the bounded-decoder
constraint belong in the exact minimization problem. Neither can be
removed by a rank calculation.

\section{Exact memory under terminal readout contrast}
\label{sec:r74-contrast}

The normal form separates the cost of a permutation module from the positivity
of its preparation and decoder. We now determine when that linear lower bound
is attained, for every finite observable group, and describe the exact memory
transitions as the terminal contrast varies. This is a change of terminal
readout. It differs from the length-dependent background damping of
Section~\ref{sec:r73-finite-certificates}.

Let $b$ be a nonconstant compactly reversible response with finite observable
group $G$, residual cardinality $N$, and real observable module $V_b$. Set
\begin{equation}\label{eq:r74-contrast}
 \pi=\frac1{|G|}\sum_{g\in G}b(g),\qquad
 b_t(g)=\pi+t\{b(g)-\pi\},\qquad 0\le t\le1.
\end{equation}
If $\pi_c=0$, nonnegativity implies $b(g,c)=0$ for every $g$; such an output
can be removed and later restored as a zero decoder column. Thus positivity
of all remaining $\pi_c$ imposes no additional hypothesis. For prepared
unitary commands, \eqref{eq:r74-contrast} is realized by the terminal
measurement $E_c(t)=tE_c+(1-t)\pi_c I$, with the preparation and commands
unchanged.

\begin{theorem}[Attainment of the module bound at small contrast]
\label{thm:r74-low-contrast}
For every finite observable group $G$ there is an explicitly positive $t_0$
such that
\begin{equation}\label{eq:r74-small-contrast}
 \kappa_G(b_t)=\nu_G(V_b),\qquad
 m(b_t)=\min\{N,1+\nu_G(V_b)\}\qquad(0<t\le t_0).
\end{equation}
For every $t>0$, the observable group, module and residual cardinality agree
with those of $b$. At $t=0$ the minimum is one label. The assertion allows
arbitrary finite groups and arbitrary finite terminal outcome sets.
\end{theorem}
\begin{proof}
Every $f\in V_b$ has a suffix-independent total mass
$\ell(f)=\sum_c f(z,c)$, because this holds for its spanning residuals.
Group averaging sends every residual to the same constant response
$\overline b(z)=\pi$. Its image on $V_b$ is therefore $\R\overline b$, and
\[
 V_b=\R\overline b\oplus W,
 \qquad W=\ker\left(\frac1{|G|}\sum_{g\in G}L_g\right).
\]
The module $W$ has no invariant vector, and $\ell(W)=0$.

Choose a $G$-set $X$ of minimum size $k=\nu_G(V_b)$ and an equivariant
injection $j:W\to\R^X$. The permutation inner product is invariant.
Since group averaging commutes with $j$, its image is orthogonal to the
whole invariant subspace, including the uniform row
$p=(1/k,\ldots,1/k)$. In particular its vectors have coordinate sum zero.
Let $F:\R^X\to W$ be orthogonal projection onto $j(W)$ followed by
$j^{-1}$. It is equivariant, $Fj=I_W$, and $F(p)=0$.

Put $v=b-\overline b$, $a=j(v)$, and
\[
 T(x,c)=F(e_x)(\varnothing,c).
\]
The rows of $T$ sum to zero, $pT=0$, and equivariance gives
$aP_gT=b(g)-\pi$. Since $a\ne0$, we may take
\[
 \epsilon=\frac{1}{2k\|a\|_\infty},\qquad
 B=\max_{x,c}\frac{|T(x,c)|}{\pi_c}>0,
 \qquad t_0=\min\{1,\epsilon/B\}.
\]
The row $\lambda=p+\epsilon a$ is a probability row with every entry at
least $1/(2k)$. Since $(t/\epsilon)B\le1$, for $0\le t\le t_0$ the matrix
\[
 D_t(x,c)=\pi_c+(t/\epsilon)T(x,c)
\]
is row-stochastic. Direct multiplication yields
\[
 \lambda P_gD_t
 =\pi+(t/\epsilon)\{pP_gT+\epsilon aP_gT\}
 =b_t(g).
\]
This proves the upper $\kappa_G(b_t)\le k$ with a bounded decoder.

The linear map
$f\mapsto tf+(1-t)\ell(f)\overline b$ is equivariant and invertible for
$t>0$: it is the identity on $\R\overline b$ and multiplication by $t$ on
$W$. It sends each residual of $b$ to the corresponding residual of $b_t$.
It preserves the group kernel and the equality pattern among residuals.
The module lower bound \eqref{eq:r73-module-bounds} now gives equality,
and Theorem~\ref{thm:r73-memory-formula} charges the pure start. At zero
contrast the decoder is constant.
\end{proof}

The construction is an equivariant positivity argument. Symmetric cone
factorizations and coordinate-symmetric lifts have an established theory
\cite{GPT13,KPT10}; ordinary nonnegative rank also has classical barycentric
perturbation results \cite{BocciCarliniRapallo11}. The assertion here is the
exact value $\nu_G(V_b)$ while preserving one common action for every command
word and the prescribed-start charge. It is not a claim that positivity
under mixing or equivariant factorization was previously unknown.

There is a useful distinction at an infinite compact observable group.
The same invertible change of residuals preserves that group for every
$t>0$, so Theorem~\ref{thm:r72-reversible} still excludes every finite
exact stochastic realization. Thus making the terminal probabilities
strictly interior need not make exact memory finite. Terminal readout
contrast, repeated channel damping, and perturbations of an entire
instrument are different operations.

\begin{theorem}[Attained exact contrast thresholds]
\label{thm:r74-contrast-thresholds}
For each integer $k\ge1$ there is $\tau_k\in[0,1]$ such that
\begin{equation}\label{eq:r74-threshold-set}
 \kappa_G(b_t)\le k\quad\Longleftrightarrow\quad 0\le t\le\tau_k.
\end{equation}
These thresholds are nondecreasing in $k$, satisfy $\tau_N=1$, and obey
\[
 k<\nu_G(V_b)\Longrightarrow\tau_k=0,
 \qquad \tau_{\nu_G(V_b)}>0.
\]
The lower memory value is retained at every threshold itself. There are
only finitely many positive-contrast transitions. If all $b(g,c)$ are real
algebraic, all thresholds are effectively computable real algebraic numbers.
When all subgroups of $G$ are normal, they are obtained by finitely many
linear programs; rational tables then have rational thresholds and rational
optimal controllers.
\end{theorem}
\begin{proof}
For a fixed action on at most $k$ points, the variables $\lambda,D,t$ lie
in a compact product of probability simplexes and $[0,1]$. The equations
$\lambda P_gD=b_t(g)$ define a closed set, so its feasible contrasts form
a compact set. If $s>0$ is feasible, the same $\lambda$ and the decoder
$(t/s)D+(1-t/s)\mathbf1\pi$ make every $0\le t\le s$ feasible.
There are only finitely many actions on at most $k$ labelled points;
their feasible sets therefore form a closed initial interval. A single
fixed point realizes $b_0$, including when no positive contrast is feasible.

The residual machine gives $\kappa_G(b_t)\le N$; the remaining inequalities
follow from Theorem~\ref{thm:r74-low-contrast} and its module lower bound.
The exact memory is recovered from
$m(b_t)=\min\{N,1+\kappa_G(b_t)\}$ for $t>0$.
For algebraic data the equations for each action are finite polynomial
constraints over the real algebraic numbers. Classical real quantifier
elimination computes their projected interval and endpoint, as in
Section~\ref{sec:r73-algebraic}. This is a proved finite procedure,
not a reported general quantifier-elimination execution. The final assertion
is the linear reduction and rationality proved in
Theorems~\ref{thm:r74-quotient-lp} and~\ref{thm:r74-quotient-dual} below.
\end{proof}

\subsection{A sharp transition for the retained prepared-qubit experiment}
The bounded-decoder obstruction of
Proposition~\ref{prop:r73-decoder-obstruction} has a complete contrast law.
Let $u=(1,-1,0)$ and
\[
 \beta_t(g)=\frac12+\frac t4uP_gu^{\mathsf T},\qquad g\in S_3,
\]
and let $b_t=(\beta_t,1-\beta_t)$. This is the same prepared-qubit
experiment with terminal effect $t|+\rangle\langle+|+(1-t)I/2$.

\begin{proposition}[An exact qubit readout threshold]
\label{prop:r74-s3-contrast}
For $0<t\le1$, this response has $d=3$, $N=6$, and
\[
 \kappa_{S_3}(b_t)=
 \begin{cases}3,&0<t\le2/3,\\6,&2/3<t\le1,\end{cases}
 \qquad
 m(b_t)=
 \begin{cases}4,&0<t\le2/3,\\6,&2/3<t\le1.\end{cases}
\]
At $t=0$ its memory is one.
\end{proposition}
\begin{proof}
The group, rank and residual count are unchanged for $t>0$ by
Theorem~\ref{thm:r74-low-contrast}. A faithful $S_3$-set of size at most
five has exactly one natural three-point orbit and otherwise only fixed
points or one two-point sign orbit, as in the proof of
Proposition~\ref{prop:r73-decoder-obstruction}. Let $A$ be the initial
mass on its natural orbit. Write its unnormalized initial row and scalar
decoder as
\[
 p=(A/3)\mathbf1^{\mathsf T}+r,
 \qquad D=q\mathbf1+s,
 \qquad r\mathbf1=0,\quad\mathbf1^{\mathsf T}s=0.
\]
The contribution is $Aq+rP_gs$. All remaining orbits contribute a linear
combination of the constant and sign functions. The functions $rP_gs$ and
$uP_gu^{\mathsf T}$ are orthogonal to both: their group sums vanish because
$|S_3|^{-1}\sum_gP_g=\mathbf1\mathbf1^{\mathsf T}/3$, and
$\sum_g\operatorname{sgn}(g)P_g=0$ entrywise. Consequently equality of the
full responses implies
\begin{equation}\label{eq:r74-s3-coefficient}
 rP_gs=(t/4)uP_gu^{\mathsf T}\qquad(g\in S_3).
\end{equation}

Here no unproved matrix-coefficient injectivity is needed. On the row plane
$\mathbf1^\perp$, use the basis $u,(1,1,-2)$. The matrices of the identity,
$(12)$, $(23)$ and the chronological cycle $(123)$ are respectively
\[
 I_2,\quad
 \begin{pmatrix}-1&0\\0&1\end{pmatrix},\quad
 \frac12\begin{pmatrix}1&1\\3&-1\end{pmatrix},\quad
 \frac12\begin{pmatrix}-1&1\\-3&-1\end{pmatrix}.
\]
They span all real two-by-two matrices: the first two span the diagonal
ones, and the latter two have independent off-diagonal parts. Therefore
\eqref{eq:r74-s3-coefficient} holds with every endomorphism of this plane
in place of $P_g$. Equality of the associated rank-one bilinear functionals,
with $t>0$, forces $r=au$, $s=bu^{\mathsf T}$ and $ab=t/4$.
Nonnegativity of $p$ gives $|a|\le A/3$, while the bounded decoder gives
$|b|\le\min(q,1-q)\le1/2$. Hence
\[
 t/4=|ab|\le A/6\le1/6.
\]
No core of size at most five is possible above $t=2/3$.

For $0<t\le2/3$ use the natural three-point action with
$\lambda=(2/3,0,1/3)$ and scalar decoder
$D=\tfrac12\mathbf1+(3t/4)u^{\mathsf T}$. Its entries lie in $[0,1]$
and direct multiplication gives $\beta_t$; the rank bound proves
$\kappa=3$. Above the threshold the six-residual machine supplies the
matching upper. The pure-start formula gives the stated minima and the
constant case is immediate.
\end{proof}

\section{Quotient decompositions, linear duality and sharp transitions}
\label{sec:r74-quotient-memory}

For groups whose subgroups are normal, the probability optimization in the
normal form has an exact linear reduction. This includes every finite abelian
group and allows arbitrary probability tables and terminal outcome sets.
Neither a special phase choice nor a positive-amplitude trigonometric ansatz
is imposed. The reduction uses normal stabilizers; it is not an unrestricted
claim that equivariant nonnegative factorization is linear programming.

\begin{lemma}[Absorption of within-orbit preparation]
\label{lem:r74-normal-orbits}
Suppose the stabilizers in a finite $G$-set are normal. Every probability
factorization on it can be replaced, without increasing its size, by one
having at most one orbit $H\backslash G$ for each normal subgroup $H$ and
whose initial row is concentrated at the identity coset within each orbit.
Only the mixture between different orbits is required.
\end{lemma}
\begin{proof}
On an orbit of positive initial mass, normalize its local preparation and
write its response as
\[
 q(g,c)=\sum_{Hx\in H\backslash G}\lambda(Hx)D(Hxg,c).
\]
For $h\in H$, normality gives $Hxhg=Hxg$, so $q(hg,c)=q(g,c)$.
Thus $q$ itself is a probability decoder on $H\backslash G$. Replacing
the preparation by the identity coset and the decoder by $q$ leaves every
command-word response unchanged. Orbits with the same $H$ can be merged
by averaging their decoders with their initial masses as weights.
Delete zero-mass orbits. These operations weakly decrease the size.
\end{proof}

Assume henceforth, until the Fourier specialization, that every subgroup of
$G$ is normal. Such groups are called Dedekind groups. For a nonempty set
$\mathcal H$ of distinct subgroups, put
$c(\mathcal H)=\sum_{H\in\mathcal H}[G:H]$.

\begin{theorem}[Exact quotient linear programs]
\label{thm:r74-quotient-lp}
The invariant $\kappa_G(b)$ is the minimum $c(\mathcal H)$ for which there
are $p_H\ge0$ and functions
$f_H:H\backslash G\times C\to\R_{\ge0}$ satisfying
\begin{equation}\label{eq:r74-quotient-lp}
 \sum_cf_H(q,c)=p_H\quad(q\in H\backslash G),\qquad
 \sum_Hp_H=1,\qquad
 \sum_H f_H(Hg,c)=b(g,c)\quad(g\in G,c\in C).
\end{equation}
The residual cardinality is $N=|G|$, and
$m(b)=\min\{|G|,1+\kappa_G(b)\}$.
For a rational response table a minimum controller can be chosen rational.
More generally, if the table lies in a specified real algebraic number field,
the controller can be chosen in that same field.
\end{theorem}
\begin{proof}
Apply Lemma~\ref{lem:r74-normal-orbits} to a minimum probability
factorization. Its orbit contributions give \eqref{eq:r74-quotient-lp}.
Conversely, if $p_H>0$, put $D_H=f_H/p_H$, place mass $p_H$ at its identity
coset, and use the right quotient action. If $p_H=0$, all its entries vanish
and the orbit is discarded. This constructs the whole response with at
most $c(\mathcal H)$ points. The regular orbit $H=\{1\}$ is always
available, and there are finitely many subgroup families.

The stabilizer $H_0$ of the initial residual is normal. If $h\in H_0$
and $g\in G$, moving $h$ past $g$ using normality shows that $h$ fixes
the residual after $g$. It therefore fixes their span and is the identity
in the minimal observable group. Thus $H_0=\{1\}$ and orbit-stabilizer
gives $N=|G|$. The pure-start assertion follows from
Theorem~\ref{thm:r73-memory-formula}.

For each family the displayed feasible set is a bounded closed polyhedron:
$0\le f_H(q,c)\le p_H\le1$. A feasible vertex is obtained by solving a
linear system over the coefficient field. It is rational for rational data
and remains in the given real algebraic field otherwise. The construction
uses only these entries and division by positive $p_H$. If the residual
machine is smaller than the initializer construction, its decoder is the
given table and has the same arithmetic property.
\end{proof}

This theorem optimizes the common permutation action supplied by the inherited
normal form. The relation between symmetry and positive factorizations is
classical \cite{GPT13,KPT10}, and positivity can cost more than a factorization
inside the target linear space \cite{GillisGlineur12}. The additional step
here is the absorption of preparation on normal quotient orbits, which
removes the bilinear unknowns without changing the counted resource.

\subsection{An explicit payoff dual}
Let $\mathcal P_H$ be the polytope of all probability tables on
$H\backslash G$, pulled back to $G$. For a real table $y$ on $G\times C$,
use the unnormalized pairing $\langle y,b\rangle=\sum_{g,c}y(g,c)b(g,c)$
and set
\begin{equation}\label{eq:r74-dual-support}
 h_H(y)=\sum_{q\in H\backslash G}\max_{c\in C}\sum_{g\in q}y(g,c),
 \qquad h_{\mathcal H}(y)=\max_{H\in\mathcal H}h_H(y).
\end{equation}

\begin{theorem}[Payoff certificates and exact contrast endpoints]
\label{thm:r74-quotient-dual}
The family $\mathcal H$ is feasible in \eqref{eq:r74-quotient-lp} if and
only if
\begin{equation}\label{eq:r74-payoff-test}
 \langle y,b\rangle\le h_{\mathcal H}(y)
 \quad\hbox{for every real table }y.
\end{equation}
If infeasible, it has a strict rational witness $y$. A witness of margin
$\Delta=\langle y,b\rangle-h_{\mathcal H}(y)>0$ still excludes that family
for any table $\widetilde b$ with
$\|\widetilde b-b\|_\infty<\Delta/\|y\|_1$.

For the contrast line \eqref{eq:r74-contrast}, put $v=b-\pi$. Its endpoint
for this family is
\begin{equation}\label{eq:r74-threshold-dual}
 \tau_{\mathcal H}
 =\min\left\{1,
  \min_{\langle y,v\rangle=1}
       \bigl(h_{\mathcal H}(y)-\langle y,\pi\rangle\bigr)\right\}.
\end{equation}
Both minima are attained. Equivalently, maximize $t$ in $[0,1]$ subject to
\eqref{eq:r74-quotient-lp} with right side $\pi+tv$.
The endpoint for budget $k$ is the maximum of these linear-program values
over $c(\mathcal H)\le k$. Rational tables give rational endpoints and
optimal controllers; no polynomial bound on the number of families is claimed.
\end{theorem}
\begin{proof}
The feasible response set for a fixed family is exactly
$\mathcal P_{\mathcal H}=\operatorname{conv}\bigcup_H\mathcal P_H$.
Each $\mathcal P_H$ is a product of probability simplexes. Maximizing a
payoff independently on each coset gives its support function $h_H$;
taking a convex hull of the union gives $h_{\mathcal H}$.
Membership proves \eqref{eq:r74-payoff-test}. If $b$ is outside this compact
convex set, let $q$ be a closest point. The first variation along the
segment from $q$ to any point of the set gives
$\langle b-q,x-q\rangle\le0$. With $y=b-q$ this yields a strictly
positive gap of at least $\|b-q\|_2^2$. Approximate $y$ by a rational
table sufficiently closely to preserve that gap; the support function is
a finite maximum of linear functions and is continuous. The perturbation
assertion is the inequality
$|\langle y,\widetilde b-b\rangle|\le\|y\|_1
\|\widetilde b-b\|_\infty$.

The constant table $\pi$ belongs to every $\mathcal P_H$. Along its ray
$\pi+tv$, a support inequality with $\langle y,v\rangle\le0$ imposes
no restriction for $t\ge0$; every other inequality can be normalized by
$\langle y,v\rangle=1$. This gives \eqref{eq:r74-threshold-dual}.
Attainment requires no strict-feasibility assumption. If $v$ lies outside
the linear space parallel to the affine hull of $\mathcal P_{\mathcal H}$,
an affine-hull equation gives an attained inner minimum zero. Otherwise
the bounded polytope has a finite facet description in that affine hull;
the first boundary reached along the ray activates a facet with positive
directional coefficient. Its normalized supporting functional attains
the inner minimum. The polytope is bounded and $v\ne0$, so the ray endpoint
is finite.

One explicit linear formulation of the inner dual introduces $z_{Hq}$
and $s$ with
\[
 z_{Hq}\ge\sum_{g\in q}y(g,c)\quad(H,q,c),\qquad
 s\ge\sum_qz_{Hq}\quad(H\in\mathcal H),\qquad
 \langle y,v\rangle=1,
\]
and minimizes $s-\langle y,\pi\rangle$. The primal formulation is obtained
by replacing the right side of \eqref{eq:r74-quotient-lp} by $\pi+tv$.
It is bounded and contains $t=0$. Vertex arithmetic proves the endpoint
and controller assertions, and the finite maximum over admissible families
gives the budget endpoint. These are exact linear certificates in the
specified finite problem, not a formal proof-assistant verification.
\end{proof}

\subsection{Fourier support and a sharp positivity transition}
For an abelian $G$, let $S$ be the complex characters whose vector Fourier
coefficient in $b$ is nonzero. It contains the trivial character and is
closed under conjugation.

\begin{corollary}[The abelian module cover]
\label{cor:r74-abelian-cover}
The number $\nu_G(V_b)$ is the minimum
$\sum_{H\in\mathcal H}[G:H]$ over distinct subgroup families such that
each character in $S$ is trivial on at least one $H\in\mathcal H$.
This lower bound is the exact value of $\kappa_G(b_t)$ for all sufficiently
small positive contrasts. It need not be attained at full contrast.
\end{corollary}
\begin{proof}
Character averaging projects the cyclic vector $b$ onto each of its
nonzero Fourier modes. Thus its complex cyclic module contains each
active character once; conjugate pairs give the corresponding real
two-dimensional summands. The permutation module on $H\backslash G$
contains exactly the characters trivial on $H$, each once. Containment
of $V_b$ is therefore exactly the stated cover. Repeated subgroups add no
new characters and may be deleted. Apply
Theorem~\ref{thm:r74-low-contrast}.
\end{proof}

The following example computes a complete rational threshold and an explicit
dual obstruction. Let $G=(\mathbb F_2)^4$, with the four coordinate flips as
commands, and let $\chi_1,\ldots,\chi_4$ be its coordinate sign characters.
Put $\chi_5=\chi_1\chi_2\chi_3\chi_4$ and
\begin{equation}\label{eq:r74-five-character-response}
 \beta_t(g)=\frac38+\frac t8\sum_{i=1}^5\chi_i(g),\qquad
 b_t(g)=(\beta_t(g),1-\beta_t(g)),\qquad 0\le t\le1.
\end{equation}

\begin{proposition}[A sharp transition invisible to the observable module]
\label{prop:r74-circuit-transition}
For every $t>0$, \eqref{eq:r74-five-character-response} has observable
group $(\mathbb F_2)^4$, dimension six, residual cardinality sixteen, and
$\nu_G(V_b)=10$. Its exact positive core and pure-start memory are
\[
 \kappa_G(b_t)=
 \begin{cases}10,&0<t\le3/5,\\12,&3/5<t\le1,\end{cases}
 \qquad
 m(b_t)=
 \begin{cases}11,&0<t\le3/5,\\13,&3/5<t\le1.\end{cases}
\]
At zero contrast the minimum is one. In particular a twelve-label minimum
does not occur in this family.
\end{proposition}
\begin{proof}
The five signs have product one, so their sum is $5,1$, or $-3$.
Consequently $\beta_1$ takes the values $1,1/2,0$, and every $\beta_t$
is a valid convex combination with the constant $3/8$.
Any four of the five characters are independent over $\mathbb F_2$.
For $t>0$ all five character coefficients and the constant mode are
nonzero, giving dimension six and trivial joint kernel. Thus $G$ is the
observable group and $N=16$.

For a quotient orbit write $W=H^\perp$ and $r=\dim W$; its size is $2^r$.
A proper $W$ contains at most $r$ of the five active nontrivial characters,
and $2^r\ge2r$ for $r=1,2,3$. An orbit with $r=4$ costs sixteen.
Covering the five modes therefore costs at least ten, and the five
two-point quotients attain this module bound.

Suppose that a positive core has at most eleven points. A size-eight
orbit covers at most three active characters; the remaining budget at
most three covers at most one more. Size sixteen is unavailable.
The nontrivial orbits thus have sizes two or four. Their total cost
must be ten and their active-character blocks partition the five modes
into singletons or pairs: each dimension must carry an active character
and no active character can be covered twice. At most one fixed point
can remain. Products belonging to different pairs are distinct and
inactive, because equality would give a nontrivial relation involving
at most four of the five circuit characters.

For this family of blocks let $\mathcal B_2$ denote its pairs and define
\[
 z(g)=1-\sum_{i=1}^5\chi_i(g)
       +\sum_{\{i,j\}\in\mathcal B_2}\chi_i(g)\chi_j(g).
\]
Averaging $z$ over the stabilizer of a singleton orbit gives $1-\chi_i$;
over a pair orbit it gives $(1-\chi_i)(1-\chi_j)$; over a fixed orbit
it gives $1$. The absence of any relation of length at most four ensures
that no other term survives these averages. Each resulting function is
nonnegative. Hence any nonnegative quotient contributions satisfy
\[
 \frac1{16}\sum_gz(g)\beta(g)
 =\sum_H\frac1{16}\sum_gz(g)f_H(Hg,1)\ge0.
\]
For the target \eqref{eq:r74-five-character-response}, character
orthogonality gives the value $(3-5t)/8$. Thus every core of size at
most eleven requires $t\le3/5$. This is also an explicit instance of
Theorem~\ref{thm:r74-quotient-dual}: take
$y(g,1)=-z(g)/16$, $y(g,0)=0$. Its family support value is zero and its
strict margin above the threshold is $(5t-3)/8$.

For $0<t\le3/5$, use five two-point quotient orbits, putting initial
mass $1/5$ at the plus point of each. The first-coordinate decoder on
orbit $i$ is
\[
 D_i(g,1)=\frac38+\frac{5t}{8}\chi_i(g).
\]
Its values lie in $[0,1]$, and the weighted sum is $\beta_t$. This
attains the ten-point core.

For $t=1$, use the four-point quotient on $(\chi_1,\chi_2)$ and the
eight-point quotient on $(\chi_3,\chi_4,\chi_5)$, with initial mass
$1/2$ at the all-plus point of each. Their decoders are respectively
the indicator of two plus signs and the indicator of at least two plus
signs. For signs $a,b,c$, the latter is
$1/2+(a+b+c-abc)/4$. The weighted sum of the two decoders is
\[
 \frac{1+\chi_1+\chi_2+\chi_1\chi_2}{8}
 +\frac14+
   \frac{\chi_3+\chi_4+\chi_5-\chi_3\chi_4\chi_5}{8}
 =\beta_1,
\]
because $\chi_1\chi_2=\chi_3\chi_4\chi_5$. Replace each decoder $D$ by
$(1-t)3/8+tD$ for other contrasts. This proves the twelve-point upper
everywhere, matching the strict lower above $3/5$.
The pure-start formula now gives eleven and thirteen, since $N=16$.
\end{proof}

\begin{corollary}[A robust finite-table obstruction]
\label{cor:r74-circuit-robust}
Fix $3/5<t\le1$ and a binary probability table $\widetilde\beta$ on the
same group $G=(\mathbb F_2)^4$. Supply it under the same four generator
commands, so the whole response is $w\mapsto\widetilde\beta(g(w))$.
If
\begin{equation}\label{eq:r74-circuit-robust}
 \max_{g\in G}|\widetilde\beta(g)-\beta_t(g)|
 <\frac{5t-3}{18},
\end{equation}
its exact positive core has at least twelve points and its pure-start
memory is at least thirteen. This is a sufficient lower certificate on
finite group tables. It does not assert a sharp neighbourhood radius or a
matching thirteen-label upper for every perturbed table.
\end{corollary}
\begin{proof}
Write $\eta$ for the left side of \eqref{eq:r74-circuit-robust}.
Every normalized Fourier coefficient changes by at most $\eta$.
Since $(5t-3)/18<t/8$ for $t\le1$, all five prescribed coefficients
remain nonzero. Their joint kernel is trivial, so the perturbed table
still has observable group $G$ and sixteen residuals. An action of size
at most eleven must therefore cover those same five modes, and the
preceding proof gives its singleton-and-pair partition even if additional
Fourier coefficients are present.

There are zero, one, or two pairs in such a partition. Every permutation
of the five circuit characters extends to a linear automorphism of the
dual group, because any four are a basis and the fifth is their sum.
Thus the absolute sum of its witness $z$ depends only on the number of
pairs. Direct evaluation on the sixteen sign patterns gives the following
values and multiplicities:
\[
 \begin{array}{c|l|c}
 \text{pairs}&\text{value : multiplicity of }z
       &16^{-1}\sum_g|z(g)|\\ \hline
 0&-4:1,\ 0:10,\ 4:5&3/2\\
 1&-3:1,\ -1:6,\ 1:4,\ 3:2,\ 5:3&17/8\\
 2&-2:5,\ 0:4,\ 2:2,\ 4:4,\ 6:1&9/4
 \end{array}
\]
The normalized pairing with $\widetilde\beta$ is at most
\[
 \frac{3-5t}{8}+\frac94\eta<0.
\]
It cannot have the required nonnegative quotient decomposition. Hence
$\kappa_G(\widetilde b)\ge12$, and
$m(\widetilde b)=\min\{16,1+\kappa_G(\widetilde b)\}\ge13$.
The upper thirteen on the exact contrast line was supplied by a specific
construction and is not inferred for arbitrary perturbations.
\end{proof}

For an exact physical realization one may use the sixteen-dimensional
diagonal model: prepare the identity basis state, use regular permutation
unitaries for the four generators, and take the diagonal terminal effect
with entries $\beta_t(g)$. This is already classical; no quantum advantage
is asserted for this example. The general low-contrast theorem, the normal
quotient linear programs, and this exact transition have different scopes.

\section{The retained instrument realization theory and its boundary}
\label{sec:r74-instruments}

The group results concern terminal responses. For an outcome at every call,
one must instead preserve each command--outcome history and every subsequent
decision made from that visible history. This stronger interface is already
covered, under a sufficient uniform posterior condition, by the earlier
instrument theory of this paper. We recall its exact scope before proving
the new conditional-memory classification.

For fixed known completely positive branches $\mathcal I_{a,y}$ whose
command sums are trace preserving, the hypothesis is
\[
 \mathcal I_{a,y}(X)\succeq
   \frac\eta h\tr\bigl(\mathcal I_{a,y}(X)\bigr)I
 \quad(X\succeq0),\qquad \eta>0.
\]
It is a condition on every recorded branch, including the zero-branch
convention, and follows when depolarizing noise of weight $\eta$ is applied
after each branch. In every finite dimension $h\ge2$, and for $k$ prescribed
preparations, Theorem~\ref{thm:r68-instrument} gives at most
\[
 k+\left\lceil(4h/\eta)^{h-1}\right\rceil
\]
hidden labels, with at most $h^2$ successors per branch row. All commands use
the same physical dictionary. The complete proof in
Section~\ref{sec:interior-instruments} gives the operator identity for every
history; it preserves all adaptive stopped laws, including nontermination,
and nonnegative costs. The hidden label is not exposed to the observer.
Corollary~\ref{cor:r68-risk} consequently preserves the risk of every common
report-based observer, including constrained boundary budgets; it does not
identify hidden simulator size with optimal observer memory.

The same retained theory includes a terminating algebraic construction
(Proposition~\ref{prop:r68-algebraic}), interiority after a finite block with
the clock and intermediate reports counted (Theorem~\ref{thm:r68-block}),
and matching powers of $\eta^{-1}$ for universal state-positive
representation (Theorem~\ref{thm:r68-noise-exponent}). The latter lower bound
is for generators which are actual density matrices. It is not imported
as a lower bound against arbitrary nonnegative hidden-state representations.
The family in Section~\ref{sec:r74-mixing-obstruction} will have a lower
bound against that unrestricted class.

Invariant positive cones and their stochastic realizations are classical
\cite{Heller65,MonrasWinter14}. Okay's quantum simulation framework explicitly
includes invariant-polytope updates and adaptive composition
\cite{Okay26}. Pure-state nets and approximation of the density-matrix body
have an established convex-geometric theory \cite{AubrunSzarek17}; extension
complexity is a different resource \cite{Fawzi21}. These ingredients, and
the paper's already proved interior-instrument theorem, are not asserted
as new discoveries of this revision. Their role is to fix the sufficient
condition and the resource convention against which the new sharp examples
are compared. The next common-contraction construction gives an explicit
certificate in supplied linear-response coordinates; it is also an
application of positive-realization geometry, not a replacement for the
retained instrument theorem.

\section{Joint contraction and bounded positive realizations}
\label{sec:r74-joint-contraction}

The reversible classification has a complementary dissipative regime. The
condition below is imposed on all command products through one supplied norm,
and the construction gives one machine for all words. A separate realization
for each command would not establish that conclusion. Invariant-polytope
positive realization is classical \cite{Heller65,MonrasWinter14}; the point of
the statement is an explicit sufficient certificate with bounded probability
decoders, a prescribed pure start, and a stated number of live labels.

Let $A$ and $C$ be nonempty finite sets, $a=|A|$, and let
\begin{equation}\label{eq:r74-contracting-response}
 b(w,c)=p_c+\ell_c(x_w),\qquad
 x_{wa}=F_ax_w,\qquad x_\varnothing=x_0\in\R^n,
\end{equation}
where $F_a$ are linear, $p\in\Delta(C)$, $\sum_c\ell_c=0$, and the displayed
values are probabilities for every word. For a norm $\|\cdot\|$ write
$\|\ell\|_*$ for its dual norm. Suppose we are given $0<\rho<1$ and $r>0$
such that
\begin{equation}\label{eq:r74-contraction-certificate}
 \|F_ax\|\le\rho\|x\|\quad(a\in A),\qquad
 r\|\ell_c\|_*\le p_c\quad(c\in C).
\end{equation}
In particular, strictly positive $p_c$ give such an $r$ for every finite
collection of output functionals. Zero components are allowed only when the
corresponding functional vanishes. The second condition certifies that the
entire ball $rB$, not only the actual orbit, has legal decoder rows.

\begin{theorem}[A common contracting realization]\label{thm:r74-contraction}
Under \eqref{eq:r74-contraction-certificate}, the response
\eqref{eq:r74-contracting-response} has an exact stochastic realization from
a prescribed pure label, valid for every command word, with at most
\begin{equation}\label{eq:r74-contraction-count}
 \sum_{j=0}^{T-1}a^j+
 \left\lceil\left(1+\frac{2}{1-\rho}\right)^n\right\rceil
 \quad\hbox{live labels},
\end{equation}
where $T\ge1$ may be any integer satisfying
$\rho^{T-1}\|x_0\|\le r$. For $x_0\ne0$ one may take
\[
 T=1+\max\left\{0,
 \left\lceil\frac{\log(\|x_0\|/r)}{\log(1/\rho)}\right\rceil\right\};
\]
for $x_0=0$, the response is constant and one label suffices. No minimality or
computational-efficiency assertion is made for this upper bound.
\end{theorem}
\begin{proof}
Put $\varepsilon=1-\rho$. Choose a maximal $\varepsilon$-separated finite
set $V$ in the unit ball $B$, where distinct selected points have distance
greater than $\varepsilon$. The open balls of radius $\varepsilon/2$ about
these points are pairwise disjoint and lie in $(1+\varepsilon/2)B$.
Lebesgue volume, which scales by the $n$th power, gives
\[
 |V|\le (1+2/\varepsilon)^n.
\]
This bound also proves that the greedy selection stops; maximality says
that $V$ is an $\varepsilon$-net of $B$. If $P=\operatorname{conv}(V)$,
then $\rho B\subseteq P\subseteq B$. Indeed, for every real linear
functional $f$, choose a maximizer $x\in B$ and a point $v\in V$ at
distance at most $\varepsilon$ from $x$. Symmetry of the unit ball gives
\[
 \max_{v\in V} f(v)\ge f(x)-\varepsilon\|f\|_*
 =(1-\varepsilon)\|f\|_*.
\]
The separating-hyperplane characterization of convex containment proves
$\rho B\subseteq P$. Consequently
\[
 F_a(rP)\subseteq\rho rB\subseteq rP
 \quad(a\in A).
\]

Use one recurrent label for each $v\in V$, representing the vector $rv$.
For every such label and command choose nonnegative barycentric coefficients
of $F_a(rv)$ in $rV$; they sum to one and define a row of $Q_a$. Give that
label decoder $p_c+\ell_c(rv)$. The decoder is nonnegative by the second
condition in \eqref{eq:r74-contraction-certificate}, and its entries sum
to one because $p$ has mass one and $\sum_c\ell_c=0$.

In addition use one transient label for each word $w$ with $|w|<T$,
representing its actual vector $x_w$. Its decoder is the actual probability
row $b(w)$, legal by hypothesis. If $|w|<T-1$, command $a$ sends it
deterministically to the label $wa$. If $|w|=T-1$, then
\[
 \|F_ax_w\|\le\rho^T\|x_0\|\le\rho r,
\]
so $F_ax_w$ is in $rP$. Choose its barycentric row in the recurrent labels.
Start at the single empty-word label. There are exactly
$\sum_{j=0}^{T-1}a^j$ such transient labels before any duplicate states
are merged. Induction on the word shows that the probability-weighted
represented vector is always $x_w$: the transition identity holds on
each individual label, and linearity preserves it under a mixture. Pairing
with all $\ell_c$ and adding $p$ gives the entire response. This proves
the count, including the pure initial label and every first command.
\end{proof}

The assumption may also be supplied as a uniform product bound. If a known
norm $\|\cdot\|_0$ and constants $K\ge1$, $0<q<1$ satisfy
\[
 \|F_w\|_0\le Kq^{|w|}\qquad(w\in A^*),
\]
choose any $q<\rho<1$ and set
\[
 \|x\|'=\sup_{w\in A^*}\rho^{-|w|}\|F_wx\|_0.
\]
This is a norm with $\|x\|_0\le\|x\|'\le K\|x\|_0$ and
$\|F_ax\|'\le\rho\|x\|'$. Hence
Theorem~\ref{thm:r74-contraction} applies with
$\|x_0\|'\le K\|x_0\|_0$ and any $r$ satisfying
$r\|\ell_c\|_{0,*}\le p_c$. This is a consequence of a supplied uniform
certificate, not a decision procedure for joint spectral radius. A common
Euclidean contraction certificate, for example, is the finite collection
$F_a^{\mathsf T}HF_a\preceq\rho^2H$ for a supplied positive definite $H$.
The actual net and barycentric rows provide finite implementation data;
their exact algebraic or sampling encodings remain distinct resources.


\section{Average mixing does not control conditional response memory}
\label{sec:r74-mixing-obstruction}

Positive realization via a common invariant polytope, including composition
of adaptive instruments, is established machinery
\cite{MonrasWinter14,Okay26}. A hypothesis forcing every normalized branch
into an interior region is stronger than fast convergence of the channel
obtained by ignoring its outcomes. The following exact example separates
these statements even when every outcome has a probability bounded away
from zero. The infinite-memory conclusion is a response-level consequence
of the inherited relative spectral obstruction; it is not a claim that
ordinary mixing theory has failed.

The phase obstruction itself has classical predecessors. The FRDN processes
discussed in \cite[Section 2.3]{Fanizza24} already exhibit
irrational-phase obstructions to finite classical realization and
phase-dependent finite minima in commensurate cases. Our assertion concerns
the simultaneous qubit, rank-three, uniform-positivity and exact
depolarization properties below, and the sharp additional charge for a
counted pure start; it does not claim priority for the phase mechanism.

Fix $\sin\theta\ne0$. For a qubit put
\[
 U_\theta=\begin{pmatrix}
 \cos(\theta/2)&-\sin(\theta/2)\\
 \sin(\theta/2)&\cos(\theta/2)
 \end{pmatrix},\qquad
 E=\begin{pmatrix}3/4&0\\0&1/4\end{pmatrix},\qquad
 \tau=|0\rangle\langle0|,
 \qquad \mathcal D(X)=\tr(X)I/2.
\]
Define three linear maps
\begin{equation}\label{eq:r74-depolarizing-instrument}
 \Psi_0(X)=\tfrac18U_\theta XU_\theta^*,\quad
 \Psi_1(X)=\tfrac14\tr(EX)\tau,\quad
 \Psi_2=\mathcal D-\Psi_0-\Psi_1.
\end{equation}

\begin{theorem}[Exact memory in a uniformly positive one-dependent family]
\label{thm:r74-depolarizing-obstruction}
The maps \eqref{eq:r74-depolarizing-instrument} form a quantum
instrument. Its unconditioned channel is the completely depolarizing map
$\mathcal D$. For every density matrix, its three outcome probabilities
belong respectively to
\[
 \{1/8\},\qquad[1/16,3/16],\qquad[11/16,13/16].
\]
Started at $I/2$, its stationary outcome process is one-dependent: two
finite blocks separated by at least one unobserved outcome are independent.
Its word-probability Hankel matrix has rank exactly three. Write
$z=e^{i\theta}$. If $z$ has finite order $q$, the least classical
hidden-state count reproducing every word probability is exactly $q$ when
a mixed initial row is free, and exactly $q+1$ when the initial label must
be pure and counted. If $z$ is not a root of unity, no finite classical
model exists, even with a freely mixed start and arbitrary real stochastic
entries. All primitive superoperators are algebraic whenever
$\cos\theta$ and $\sin\theta$ are algebraic.
\end{theorem}
\begin{proof}
Use the Choi convention
$J(\Phi)=\sum_{i,j=0}^1|i\rangle\langle j|\otimes
\Phi(|i\rangle\langle j|)$. Then
\[
 J(\mathcal D)=\tfrac12I_4,
 \qquad \|J(\Psi_0)\|_{\mathrm{op}}=\tfrac14,
 \qquad \|J(\Psi_1)\|_{\mathrm{op}}=\tfrac3{16}.
\]
The first two branch Choi matrices are positive semidefinite. Subtraction
therefore gives $J(\Psi_2)\succeq I_4/16$, proving complete positivity.
The sum is trace preserving by construction. The probability bounds follow
from $\tr\Psi_0(\sigma)=1/8$ and
$\tr\Psi_1(\sigma)=\tr(E\sigma)/4\in[1/16,3/16]$.
Conjugation by $U_\theta$ depends algebraically on $\cos\theta$ and
$\sin\theta$, which proves the last assertion.

The initial state $\sigma_*=I/2$ is fixed by $\mathcal D$, giving
stationarity. To check one-dependence directly, write $\Psi_u$ for the
chronological branch product of a block $u$. Summing over one intervening
outcome gives
\[
 \sum_{s=0}^2\tr\bigl(\Psi_v\Psi_s\Psi_u(\sigma_*)\bigr)
 =\tr\bigl(\Psi_v\mathcal D\Psi_u(\sigma_*)\bigr)
 =\tr\Psi_u(\sigma_*)\,\tr\Psi_v(\sigma_*).
\]
Further unobserved outcomes leave $\mathcal D$ unchanged. Summing over
specified block events proves the claimed independence.

After outcome $1$ the normalized state is $\tau$, and that first outcome
has probability $1/8$. Repeating outcome $0$ rotates this state with
the factor $8^{-n}$. The probability of the full word $1\,0^n\,1$ is
\begin{equation}\label{eq:r74-conditioned-oscillation}
 u_n=\frac1{64}\,8^{-n}
       \left(1+\frac12\cos(n\theta)\right).
\end{equation}
All reachable subnormalized states are real symmetric matrices, a space
of real dimension three, so the word-probability Hankel rank is at most
three. For the reverse inequality take prefixes $1\,0^i$ and suffixes
$0^j\,1$, $0\le i,j\le2$. Their Hankel submatrix is $(u_{i+j})$.
Equation~\eqref{eq:r74-conditioned-oscillation} expresses it as the
product of a three-point Vandermonde matrix, a diagonal matrix with three
nonzero coefficients, and the transpose Vandermonde matrix. It is therefore
nonsingular, proving rank three.

Suppose a finite classical model exists. Write its nonnegative outcome
matrices as $T_0,T_1,T_2$, with $\sum_sT_s$ row-stochastic, and its initial
probability row as $\mu$. Then
\[
 u_n=vT_0^nd,\qquad v=\mu T_1\ge0,\quad d=T_1\one\ge0.
\]
Here $T_0$ is row-substochastic. Delete states unreachable from the support
of $v$ and states from which no positive terminal value of $d$ is reachable.
The remaining principal submatrix $S$ still realizes $u_n$. For every
remaining $i,j$, choose a positive path from $v$ to $i$, of length $a_i$
and weight $p_i>0$, and a path from $j$ to a positive terminal value, of
length $b_j$ and weight $r_j>0$. Nonnegativity gives
\[
 p_ir_j(S^n)_{ij}\le u_{a_i+n+b_j}\le C_{ij}8^{-n}.
\]
Thus $(8S)^n$ is entrywise bounded, and every eigenvalue of $S$ has modulus
at most $1/8$.

The three exponential modes of \eqref{eq:r74-conditioned-oscillation}
are $1/8$, $z/8$, and $\overline z/8$, all distinct and with nonzero
coefficients. Cayley--Hamilton and the Vandermonde independence of these
sequences force $z/8$ to be an eigenvalue of $S$. In an upper block
triangular decomposition of $S$, it belongs to an irreducible nonnegative
diagonal block $B$. Its spectral radius is exactly $1/8$. Perron
normalization by a strictly positive eigenvector therefore makes $8B$
similar to a stochastic matrix with eigenvalue $z$. This is the elementary
normalization proved before Proposition~\ref{prop:r70-relative-spectrum}.
At a largest-modulus coordinate of a corresponding complex eigenvector,
equality in the triangle inequality forces each positive successor to have
the same modulus and value multiplied by $z$. Irreducibility propagates this
equality to every coordinate, so the relation holds on every positive edge
and every directed cycle. Following positive successors
reaches a directed cycle; its length $l\ge1$ gives $z^l=1$. This
identity proves that $z$ must be a root of unity. If its order is $q$,
every such directed cycle has length divisible by $q$, so every
realization has at least $q$ labels.

For the matching upper, suppose $z$ has order $q\ge3$. In the real Bloch
plane let $P$ be the regular $q$-gon with vertices
\[
 r_j=(\cos(2\pi j/q),\sin(2\pi j/q)),\qquad 0\le j<q,
\]
and let $\sigma_j$ be the corresponding pure qubit states. The polygon
contains the disk of radius $\cos(\pi/q)\ge1/2$. The normalized branch
$0$ permutes its vertices and branch $1$ resets to $\sigma_0=\tau$.
For an arbitrary state with Bloch vector $r$ in the unit disk, write
$h=\tr\Psi_1(\sigma)\in[1/16,3/16]$. The normalized Bloch vector after
branch $2$ is
\[
 \frac{-R_\theta r/8-h r_0}{1-1/8-h},
\]
whose norm is at most $(1/8+3/16)/(11/16)=5/11<1/2$. It therefore lies
in $P$. For each input vertex and outcome, expand the normalized output in
the same vertices and multiply the barycentric coefficients by its outcome
probability. These nonnegative rows form outcome matrices $T_s$ with
$\sum_sT_s$ row-stochastic. They reproduce each subnormalized branch
state on each label. Induction proves equality of every word probability.
The uniform mixture of the vertices represents $I/2$, giving the
$q$-label upper with a free mixed start. A stationary hidden initialization
can also be used: with $T=\sum_sT_s$, every row $\mu T^n$ represents
$\mathcal D^n(I/2)=I/2$, where $\mu$ is the uniform row. A convergent
subsequence of their Cesaro averages gives $\pi T=\pi$, and $\pi$ still
represents $I/2$, so it gives the same word probabilities on $q$ labels.
If a pure start is required, add
one label representing $I/2$ and use its branch decompositions on the first
outcome; the first outcome is generated during this counted transition.
No transition returns to that label. This gives the $q+1$ upper.

It remains to exclude a pure-start realization on exactly $q$ labels.
The preceding lower then forces the irreducible block carrying $z/8$
to use every label. Perron normalization makes $8T_0$ similar to a
stochastic matrix with primitive $q$th-root eigenvalue and exactly $q$
states. The equality argument above supplies a directed $q$-cycle using
every state. An additional positive edge would close a directed cycle of
length less than $q$, which is impossible. Thus, before normalization,
$T_0$ is a weighted directed $q$-cycle. Its successive row weights are
positive. From the prescribed starting point,
\[
 p(0^n)=8^{-n}\qquad(n\ge0).
\]
For $n=1,\ldots,q$, the left side is the product of the first $n$ row
weights encountered on this cycle. Successive ratios force every weight
to be $1/8$, so $T_0=C_q/8$ after relabelling. Next,
$p(0^n1)=8^{-n}/8$ shows that every coordinate of $d=T_1\one$ is
$1/8$. This would give
\[
 p(1\,0^n\,1)
 =e_1^{\mathsf T}T_1(C_q/8)^nd
 =8^{-n}/64,
\]
contradicting \eqref{eq:r74-conditioned-oscillation}, already at $n=0$.
The pure-start minimum is therefore $q+1$.
\end{proof}

In particular, the fully algebraic choice
\[
 U_\theta=5^{-1/2}\begin{pmatrix}2&-1\\1&2\end{pmatrix},
 \qquad e^{i\theta}=(3+4i)/5,
\]
has no finite classical realization. Its superoperator coefficients are
rational. The phase is not a root of unity, since otherwise
$z+z^{-1}=6/5$ would be a rational algebraic integer and hence an integer.
Every finite-order choice, on the other hand, gives the exact counts
$q$ and $q+1$ in a fixed physical qubit with the same uniform outcome
probability bounds. The size of an exact hidden-state realization cannot be inferred
from predictive rank, unconditional mixing time, or finite dependence
range.

Stationary one-dependent processes without a finite hidden Markov model
were established earlier, notably by the finitely dependent colorings of
Holroyd and Liggett \cite{HolroydLiggett16}. We make no first-existence
claim for that phenomenon. Here the simultaneous properties are an explicit
qubit instrument, rank-three linear prediction, uniformly positive outcome
probabilities, and an exactly depolarizing unconditioned channel, together
with a sharp finite-order memory law and a separate pure-initialization
charge.

This example satisfies a stronger average-mixing property than exponential
convergence: one unobserved step erases all prior state information. The
infinite obstruction survives in a specified observed word, whose
conditional maps retain an irrational phase. It therefore does not
contradict a finite-realization theorem requiring uniform mixing of each
normalized branch. It also explains why positive single-outcome
probabilities cannot replace that hypothesis. The example supplies a precise boundary for the constructive theorem.


\section{Exact synthesis from algebraic observable data}\label{sec:r73-algebraic}
The preceding normal form concerns existence over the real numbers. We now
give a terminating exact procedure on algebraic input, including a certificate
that no finite stochastic realization exists. The input is a supplied
finite-dimensional probability response, rather than a finite sample from an
unknown process. All algebraic numbers are specified by defining polynomials
and isolating data, with a common real number field and its chosen embedding.
The algorithms below concern exact label count and exact controller tables.
They make no polynomial running-time or bounded-bit sampling claim.

Finite matrix-group recognition is classical, and substantially faster
algorithms than the elementary enumeration below are known
\cite{BabaiBealsRockmore93,Serre07}. Reduction of probabilistic-automaton
minimization to the existential theory of the reals is also established
\cite{KieferWachter14}. Together with the inherited reversible criterion,
these ingredients already yield decidability and bounded-size minimization
on the input class considered here. The new normal form supplies the more
structured conclusion: a smallest controller can be returned as a finite
permutation core with a precisely counted initializer and a bounded
multi-output decoder. We give a complete terminating construction, including
finite certificates of the negative answer. Neither finite-group recognition
nor real-algebraic automaton minimization is asserted to be a new algorithm.

\subsection{Observable reduction and a finite stopping bound}
Write the supplied representation in chronological row convention as
\begin{equation}\label{eq:r73-algebraic-input}
 b(w)=\beta T_w B,\qquad T_{uv}=T_uT_v,
\end{equation}
where the matrices, row and decoder have entries in a real number field
$K$ of degree $e$ over $\mathbb Q$. The displayed representation is linear;
its matrices and intermediate coordinates need not be stochastic or
nonnegative. Probability validity means that the supplied $b(w)$ belongs to
$\Delta(C)$ for every word. Compact reversibility is imposed on its
\emph{minimal observable} response space, as in
Section~\ref{sec:r72-reversible}.

For clarity, minimal reduction itself is finite. Starting with $\beta$,
close its row span under right multiplication by each $T_a$. Starting with
the columns of $B$, close their span under left multiplication by the
$T_a$. Either iteration stabilizes once its dimension fails to increase,
and neither dimension exceeds the supplied dimension. A reachable row is
invisible precisely when it annihilates the second span. Quotienting the
reachable span by this kernel gives the minimal representation, with
dimension $d=\dim V_b$ and coefficients still in $K$. This is the familiar
reachable/observable linear reduction; it does not identify arbitrary
hidden stochastic rows with observable states. Since the response is
normalized and nonempty, $d\ge1$.

Set
\begin{equation}\label{eq:r73-group-budget}
 h=ed,\qquad
 B_h=|\operatorname{GL}_h(\mathbb F_3)|
     =\prod_{j=0}^{h-1}(3^h-3^j).
\end{equation}
This deliberately conservative bound has the advantage of a short complete
proof. Restriction of scalars may use a larger field than necessary; doing
so changes the computational bound, not the resulting exact label minimum.

\begin{lemma}[Classical finite-group bound]\label{lem:r73-minkowski}
A finite subgroup of $\operatorname{GL}_d(K)$ has order at most $B_{ed}$.
Consequently a monoid generated by invertible $d$-dimensional matrices over
$K$ is finite if and only if breadth-first enumeration of its distinct
positive-word products closes before producing $B_{ed}+1$ elements.
\end{lemma}
\begin{proof}
Restriction of scalars embeds the group faithfully in
$\operatorname{GL}_{ed}(\mathbb Q)$. For a finite subgroup $H$ of
$\operatorname{GL}_h(\mathbb Q)$, let
$\Lambda=\sum_{g\in H}g\mathbb Z^h$. There is a common denominator $q$
with $\mathbb Z^h\subseteq\Lambda\subseteq q^{-1}\mathbb Z^h$.
Thus $\Lambda$ is a free abelian group of rank $h$, and it is invariant
under every $g$ and $g^{-1}$. In a lattice basis, $H$ is a subgroup of
$\operatorname{GL}_h(\mathbb Z)$.

The kernel of reduction modulo $3$ has no nonidentity element of finite
order. Indeed, such an element would have a power $A$ of prime order
$\ell$. Write $A=I+3^s Z$, with $s\ge1$ and with an integer matrix $Z$
not zero modulo $3$. The binomial expansion of $A^\ell=I$, divided by
$3^s$, gives $\ell Z=0$ modulo $3$, so $\ell=3$. The expansion for this
prime, divided by $3^{s+1}$, instead gives
\[
 Z+3^sZ^2+3^{2s-1}Z^3=0,
\]
again contradicting $Z\not\equiv0\pmod3$. Reduction is therefore
injective on $H$. Counting independent ordered columns over $\mathbb F_3$
proves the bound in \eqref{eq:r73-group-budget}. This is the standard
Minkowski argument \cite{Serre07}; the proof is included to make the
termination certificate explicit.

A finite monoid of invertible matrices is a group: for every element,
two positive powers agree, and invertibility yields a positive power
equal to the identity and another equal to its inverse. A finite generated
monoid therefore has at most $B_{ed}$ elements. In a breadth-first search,
process each newly encountered product and multiply it by all generators.
If the queue empties, the finite set found is closed under all generators
and is the entire monoid. If the monoid is infinite, the search eventually
finds $B_{ed}+1$ distinct products. These are the two asserted stopping
events. All equality tests are exact tests in $K$.
\end{proof}

\begin{theorem}[Decidability and minimal algebraic synthesis]
\label{thm:r73-algebraic-synthesis}
For a supplied compactly reversible probability response with algebraic
linear representation \eqref{eq:r73-algebraic-input}, there is a terminating
algorithm that decides whether an exact finite stochastic realization
exists. If the answer is positive, the algorithm returns its least number
$m(b)$ of live labels and an exact realization with algebraic entries,
a prescribed initial label and a row-stochastic terminal decoder. It can
return a permutation core with at most one transient initializer, as in
Theorem~\ref{thm:r73-normal-form}.

A negative answer has a finite certificate consisting of $B_{ed}+1$
distinct observable command products. A positive answer has an exact
controller, the finite observable-group multiplication table, and finite
real-algebraic feasibility or infeasibility certificates for every smaller
candidate size. The negative certificate concerns the promised compactly
reversible input class; it is not a nonrealizability test for arbitrary
dissipative responses.
\end{theorem}
\begin{proof}
Perform the observable reduction and apply Lemma~\ref{lem:r73-minkowski}
to the minimal command matrices. If enumeration exceeds its bound, the
positive-word monoid and its compact group closure are infinite. Theorem
\ref{thm:r72-reversible} rules out every finite stochastic realization.
The list of distinct products, each supplied with its generating word,
is a finite independently checkable certificate of this conclusion.

Otherwise enumeration closes to a finite group $G$. Compute the distinct
reachable rows $\beta T_g$ in the minimal representation; their count $n$
is the residual cardinality. The deterministic residual machine proves
$m(b)\le n$. Theorem~\ref{thm:r73-memory-formula} expresses the answer as
a minimum over finite permutation actions, a probability row $\lambda$,
and a stochastic decoder $D$. Only the following finite possibilities can
improve or attain the upper $n$: a core of size $k\le n$ with a point
initial row, or a core of size $k\le n-1$ with a possibly mixed initial
row and one additional initializer.

For each such $k$, enumerate assignments of permutations to the elements
of $G$ and retain those satisfying
$P_1=I$ and $P_gP_h=P_{gh}$. This is a finite search through at most
$(k!)^{|G|}$ assignments, and it may of course be reduced using generators
and relations. For a retained action impose
\begin{equation}\label{eq:r73-bilinear-synthesis}
 \lambda_i\ge0,\quad\sum_i\lambda_i=1,\qquad
 D_{ic}\ge0,\quad\sum_cD_{ic}=1,\qquad
 \lambda P_gD=b(g)\quad(g\in G).
\end{equation}
For the point case impose $\lambda=e_i^{\mathsf T}$ and enumerate $i$.
For the other case it is harmless to permit point rows as well: that
candidate then pays a redundant initializer, and the separately enumerated
point case can only improve its cost. Equations
\eqref{eq:r73-bilinear-synthesis} have algebraic coefficients and degree
at most two in their real unknowns. Exact real-algebraic elimination
decides feasibility, and a nonempty such set has an algebraic sample point.
One may see the latter directly from a sign-invariant algebraic cell
decomposition: choose isolating algebraic sample coordinates successively
in its nonempty sections and sectors. Algebraic elimination is the
classical decision theorem used here \cite{Seidenberg54}; no numerical
optimizer is being substituted for it.

Search in increasing total label cost and stop at the first feasible
candidate. Termination follows from the deterministic upper. Theorem
\ref{thm:r73-memory-formula} proves that no stochastic competitor outside
the enumerated normal forms can be smaller. For a point initial row use
the permutation matrices directly. For a mixed initial row add the
counted initial state, give it decoder $b(\varnothing)=\lambda D$, and
on command $a$ send it to $\lambda P_{g_a}$ in the core. All later
transitions are permutations. Every returned entry is algebraic, and all
response identities follow from the group action and
\eqref{eq:r73-bilinear-synthesis}.

The finite word criterion of Theorem~\ref{thm:r73-finite-tests} supplies
an alternative bounded-size formulation using arbitrary stochastic
matrices, independently of the normal form. This is useful for checking
a returned implementation against the original linear representation.
Its existence does not give an efficient algorithm for minimizing
probabilistic automata, a problem with established hardness results
\cite{KieferWachter14}.
\end{proof}

\begin{corollary}[Algebraic prepared unitary experiments]
\label{cor:r73-unitary-algorithm}
Let the prepared density matrix, a finite family of unitary commands and
a finite terminal positive-operator-valued measurement have algebraic
real and imaginary entries. Whether their Born probabilities for all
command words admit an exact finite classical stochastic realization is
decidable. In the positive case a smallest exact classical realization
and its bounded decoder can be synthesized with algebraic entries.
For a protected Hilbert space of dimension $d_{\rm phys}$, the observable
linear dimension used in the algorithm is at most $d_{\rm phys}^2$.
\end{corollary}
\begin{proof}
Hermitian matrices have a real coordinate basis of size
$d_{\rm phys}^2$ consisting of diagonal matrix units and real and
imaginary symmetric matrix-unit pairs. In this basis conjugation by each
supplied unitary, preparation, and terminal evaluation have real algebraic
coefficients. This constructs \eqref{eq:r73-algebraic-input}. Positivity
and normalization of Born probabilities follow from the declared density
matrix and measurement, not from a finite test. The observable reduction
is finite linear algebra over a common real number field. Corollary
\ref{cor:r72-unitary} proves compact reversibility of this quotient, so
Theorem~\ref{thm:r73-algebraic-synthesis} applies.
\end{proof}

This corollary allows noncommuting commands, nonfaithful physical actions,
mixed preparations and arbitrary finite terminal measurements. In
particular, a physically infinite group with constant observable response
is correctly reduced to the one-label case before group enumeration. The
result concerns exact simulation of the supplied experiment. Choosing an
optimal feedback experiment, deciding validity of an arbitrary algebraic
series, approximate program size, exact sampling and workspace lower
bounds are different questions. The complete retained causal and boundary
theory keeps its own quantifiers and resource accounts.

\section{Finite certificates and stable exact memory}\label{sec:r73-finite-certificates}

The normal form gives exact realizations at reversible responses. We next prove
that their memory obstructions admit finite tests and persist under a specified
class of physical perturbations. Finite-word equivalence for linear machines is
classical \cite{Tzeng92}; the finite bound below is its direct difference-space
argument. The additional conclusions concern the positive separation from
\emph{all} smaller stochastic machines, its effective algebraic certificate,
and a matching realization under damping toward a common background response.

\subsection{Finite tests against a fixed memory budget}
Suppose a response $b:A^*\to\Delta(C)$ has a supplied real linear representation
\[
 b(w)=\beta T_w B,
 \qquad \beta\in\R^{1\times d},\quad T_a\in\R^{d\times d},
 \quad B\in\R^{d\times |C|}.
\]
Minimality of this representation is not required. For $M\ge1$ let
$\mathcal P_M$ be the compact product of the row-stochastic matrix simplexes
for $(Q_a)_{a\in A}$ and $D\in\R^{M\times |C|}$. Its initial row is the
prescribed pure row $e_1^{\mathsf T}$, as before. Set
\begin{equation}\label{eq:r73-finite-gap}
 L=d+M-1,\qquad
 \delta_M(b)=\min_{(Q,D)\in\mathcal P_M}
 \max_{\substack{|w|\le L\\c\in C}}
 \left|e_1^{\mathsf T}Q_wD_c-b(w,c)\right|.
\end{equation}

\begin{theorem}[Finite obstruction and effective separation]\label{thm:r73-finite-tests}
An $M$-label stochastic machine realizes $b$ on every word if and only if it
does so on the words of length at most $d+M-1$. Consequently
$\delta_M(b)>0$ if and only if no exact $M$-label realization exists.
If another supplied response $\widetilde b$ satisfies
\[
 \max_{|w|\le L,\ c\in C}|\widetilde b(w,c)-b(w,c)|\le\eta,
\]
then every $M$-label machine has error at least $\delta_M(b)-\eta$ on these
same probes. This last assertion imposes no finite-rank hypothesis on
$\widetilde b$.

If the given entries of $\beta,T_a,B$ are real algebraic numbers,
$\delta_M(b)$ is an effectively computable real algebraic number, and a
minimizing machine with real algebraic entries can be computed. When it is
positive, a positive rational number strictly below it and a real-algebraic
quantifier-elimination certificate of that strict lower bound can be computed.
No polynomial-time or practical-size bound is asserted.
\end{theorem}
\begin{proof}
For a fixed candidate machine use the block-diagonal matrices
$S_a=\operatorname{diag}(T_a,Q_a)$, the row
$u=(\beta,-e_1^{\mathsf T})$, and the output columns
$v_c=(B_c,D_c)^{\mathsf T}$. The difference on $w$ is $uS_wv_c$.
Let $R_j$ be the span of $uS_w$ for $|w|\le j$. These spaces increase;
if $R_j=R_{j+1}$, the common space is invariant under every $S_a$ and all
later spaces coincide. Since $u\ne0$, the sequence starts in dimension one
inside a space of dimension $d+M$ and must have stabilized by $j=d+M-1$.
Vanishing on the indicated finite set therefore implies vanishing on all
words. This is the same linear-equivalence principle as the inherited
finite-realization test.

The maximum in \eqref{eq:r73-finite-gap} is a continuous function on the
compact set $\mathcal P_M$, so its minimum is attained. It is zero precisely
when one machine passes every finite probe, hence, by the first assertion,
realizes every word. The perturbation claim follows pointwise from the triangle
inequality and then from the minimum defining $\delta_M(b)$.

For algebraic input, every probe value is algebraic. The constraints defining
row-stochasticity and the inequalities between the finitely many polynomial
errors and a variable $t$ are polynomial constraints over the real algebraic
numbers. Real quantifier elimination decides their feasibility and computes
the least feasible $t$ together with an algebraic witness. A singleton
semialgebraic set over the real algebraic numbers consists of an algebraic
number. If $\delta_M(b)>0$, effective isolation supplies a positive rational
$\sigma<\delta_M(b)$; infeasibility with all errors at most $\sigma$ is a
finite real-algebraic certificate. This uses the classical decision procedure
\cite{Seidenberg54}; it does not assert that this procedure was executed by
the finite diagnostics accompanying this article.
\end{proof}

For the prepared quantum response, these finite probes have an explicit
physical perturbation bound. Suppose two density matrices differ in trace
norm by at most $h_\rho$, corresponding channels differ in diamond norm by
at most $h_\Phi$, and corresponding terminal effects differ in operator norm
by at most $h_E$. All channels are trace preserving and completely positive,
and both lists of effects are valid measurements. Contractivity and telescoping
give $\|\widetilde\Phi_w(\widetilde\rho)-\Phi_w(\rho)\|_1
\le h_\rho+|w|h_\Phi$. The state difference has trace zero; subtracting
$I/2$ from a terminal effect, whose spectrum lies in $[0,1]$, gives
\begin{equation}\label{eq:r73-physical-probes}
 \max_{|w|\le L,c}|\widetilde b(w,c)-b(w,c)|
 \le h_E+\tfrac12(h_\rho+Lh_\Phi).
\end{equation}
Thus a strict smaller-memory obstruction survives a specified open
neighbourhood of physical primitives. This supplies a lower bound there;
it does not assert that an arbitrary perturbation has a realization of the
same size.

\subsection{A matching upper under a common background}
Let $b$ be compactly reversible, nonconstant, and have finite stochastic
memory $m=m(b)$. Let $G$ be its finite observable group, with the chronological
convention of the normal-form theorem, and put
\[
 \overline b=|G|^{-1}\sum_{g\in G} b(g),\qquad
 C_b=\max_{g\in G,c\in C}|b(g,c)-\overline b(c)|>0.
\]
Choose a damping factor $\alpha_a\in[0,1]$ for each command, put
$\beta_w=\prod_{j=1}^{|w|}\alpha_{w_j}$, with $\beta_\varnothing=1$, and define
\begin{equation}\label{eq:r73-common-background}
 b_{\boldsymbol\alpha}(w)
   =\beta_w b(w)+(1-\beta_w)\overline b.
\end{equation}

\begin{theorem}[Stable exact memory under background damping]\label{thm:r73-memory-plateau}
The response \eqref{eq:r73-common-background} has an exact realization with
$m$ labels for every choice of the damping factors. Put $L=d+m-2$, where
$d$ is any supplied representation dimension of $b$, and let
$\epsilon=\max_a(1-\alpha_a)$. Then
\begin{equation}\label{eq:r73-memory-plateau}
 \epsilon<\min\left\{1,\frac{\delta_{m-1}(b)}{LC_b}\right\}
 \quad\Longrightarrow\quad m(b_{\boldsymbol\alpha})=m.
\end{equation}
For algebraic input satisfying the hypotheses, a positive algebraic cutoff,
or a smaller positive rational cutoff, is effectively computable.
For a constant response every damping leaves that response unchanged and its
minimum is one label; the expression involving $m-1$ is not used.
\end{theorem}
\begin{proof}
Take a minimal realization in the normal form proved above. On its recurrent
labels each $Q_a$ is a permutation, and an extra initializer is present only
when the initial distribution on these labels is mixed. Average that
distribution over $G$, and call the resulting row $\pi$, extended by zero
on the initializer when present. It obeys
\[
 \pi Q_a=\pi,\qquad \pi D=\overline b.
\]
Indeed, the group average is invariant under each permutation. The initializer
has no incoming transition, and its first-input row is the original initial
distribution followed by the commanded permutation, so the same identities
hold on the entire machine. Set $J=\mathbf1\pi$. Then
$J^2=J$, $Q_aJ=JQ_a=J$, and
\[
 Q_{a,\boldsymbol\alpha}=\alpha_aQ_a+(1-\alpha_a)J
\]
is row-stochastic on the same labels. Multiplication gives, for every nonempty
word and also trivially at the empty word,
\[
 Q_{w,\boldsymbol\alpha}=\beta_wQ_w+(1-\beta_w)J.
\]
Reading with the same decoder proves the upper.

By minimality $\delta_{m-1}(b)>0$. For $|w|\le L$, the elementary product
bound $1-\beta_w\le |w|\epsilon$ gives an error from $b$ of at most
$LC_b\epsilon$. Theorem~\ref{thm:r73-finite-tests} rules out every
$(m-1)$-label realization when this is smaller than $\delta_{m-1}(b)$.
Machines with fewer labels can be padded by unreachable rows, so they too are
excluded. This and the upper prove equality. A nonconstant response has
$m\ge2$ and $d\ge2$, hence $L>0$. The effective assertion follows from the
finite group enumeration and algebraic synthesis theorem, together with
Theorem~\ref{thm:r73-finite-tests}. Finite maxima and averages over the
computed group preserve algebraicity.
\end{proof}

There is a physical realization of the common background for prepared unitary
commands. Let $\mathcal K$ be the compact closure of their physical conjugation
group. The compact convex hull of $\{\Phi(\rho):\Phi\in\mathcal K\}$ has a
unique element $\rho_*$ of least Hilbert--Schmidt norm: strict convexity gives
uniqueness. Every conjugation in $\mathcal K$ preserves that set and norm, so
$\rho_*$ is invariant. Its observable image belongs to the convex hull of the
finite orbit of $b$ and is invariant. Averaging any element of that orbit
over $G$ gives $\overline b$; applying this averaging projection to an
invariant element in its convex hull shows that its observable image is
exactly $\overline b$. Therefore the physical channels
\[
 \widetilde\Phi_a(X)=\alpha_aU_aXU_a^*
                   +(1-\alpha_a)\operatorname{tr}(X)\rho_*
\]
produce \eqref{eq:r73-common-background} under every command word. In general
$\rho_*$ need not be maximally mixed. The theorem concerns damping toward
this common invariant background, not every depolarizing or noisy channel.

\subsection{Periodic plateaux and an explicit rational certificate}
For the periodic mixture of Section~\ref{sec:r72-periodic}, let the distinct
nonreal frequencies have positive amplitudes $\gamma_j$, total
$\Gamma=\sum_j\gamma_j$, and let $M_*$ be its exact partition-and-anchor
minimum. Its dimension is $d=2r+1$ and $C_b\le\Gamma/2$. Hence a common
damping $\alpha=1-\epsilon$ retains exactly $M_*$ labels whenever
\begin{equation}\label{eq:r73-periodic-cutoff}
 0\le\epsilon<
 \min\left\{1,\frac{2\delta_{M_*-1}(b)}{\Gamma(2r+M_*-1)}\right\}.
\end{equation}
This is an effective finite certificate for each algebraic periodic input.
It closes the previously existential plateau cutoff for this class, without
claiming a sharp transition point or a fixed-memory risk asymptotic.

For clarity, one certificate can be calculated without a general
quantifier-elimination run. Consider the unary response
\begin{equation}\label{eq:r73-period-four}
 b_\alpha(n)=\tfrac12\{1+\alpha^n\cos(\pi n/2)\},\qquad n\ge0.
\end{equation}

\begin{proposition}[A certified four-label interval]\label{prop:r73-period-four}
For the undamped sequence $b=b_1$ and the dimension $d=3$,
$\delta_3(b)>1/36$. In particular the exact stochastic memory of
\eqref{eq:r73-period-four} is four for $89/90\le\alpha\le1$.
This interval is a certified sufficient interval, not a claimed sharp one.
\end{proposition}
\begin{proof}
Let $Q$ be any row-stochastic $3\times3$ matrix, put $t=\operatorname{tr}Q$
and $h=\det Q$, and write $y_n=e_1^{\mathsf T}Q^nD$ for any scalar decoder.
Since $Q\mathbf1=\mathbf1$, its characteristic polynomial is
\[
 z^3-tz^2+(t+h-1)z-h.
\]
Thus $y_{n+3}-ty_{n+2}+(t+h-1)y_{n+1}-hy_n=0$. We have $0\le t\le3$,
$|h|\le1$, and $|t+h-1|\le3$. The determinant bound follows directly by
bounding its absolute permutation sum by the product of the row sums.
If the first five values of $y$ are within $\eta$ of
$(1,1/2,0,1/2,1)$, substitution at $n=0,1$ yields
\[
 |t-h|\le16\eta,\qquad |2-t-h|\le16\eta,
 \qquad |t-1|,|h-1|\le16\eta.
\]
Nonnegativity and the row sums imply a further exact identity. For
$a_i=1-Q_{ii}$ it is
\begin{equation}\label{eq:r73-three-state-certificate}
 t^2-3t+3-3h
 =\tfrac12\sum_{i<j}(a_i-a_j)^2
       +3\sum_{i<j}Q_{ij}Q_{ji}\ge0.
\end{equation}
To verify it, the second characteristic coefficient is both $t+h-1$ and
$\sum_{i<j}(Q_{ii}Q_{jj}-Q_{ij}Q_{ji})$; substitute
$Q_{ii}=1-a_i$ and simplify. On the other hand, putting $u=t-1$ and $v=h-1$
gives
\[
 t^2-3t+3-3h=-2-u+u^2-3v
       \le-2+64\eta+256\eta^2.
\]
For $\eta\le1/36$ the right side is at most $-2/81<0$, a contradiction.
The compact minimum on the first five probes is therefore strictly greater
than $1/36$, and so is \eqref{eq:r73-finite-gap}, which here uses the six
probes $0\le n\le5$.

The undamped response has dimension three and a four-cycle realization with
decoder $(1,1/2,0,1/2)$. The same polynomial certificate excludes three labels,
so its minimum is four. Damping that cycle toward the uniform stationary row
realizes \eqref{eq:r73-period-four} with four labels. Its difference from the
undamped sequence on $0\le n\le5$ is at most
$\tfrac52(1-\alpha)$. For $\alpha\ge89/90$ this is at most $1/36$, strictly
below $\delta_3(b)$. The finite obstruction proves the lower and completes
the proof.
\end{proof}

The certificate uses a stochastic-matrix inequality rather than a sample
search over three-label machines. Its accompanying rational diagnostics check
the displayed polynomial identity, the constructed all-word damping
identities on explicit finite models, and exact finite recurrence consequences.
They do not replace the proofs, compute all algebraic gaps, or give an
independent human or editorial assessment.

\section{Relation to the retained exact resource laws}
\label{sec:r73-consequences}
The finite-group normal form has a cyclic specialization, but it does not
replace the independent spectral and arithmetic arguments for responses
outside the reversible class. We record the precise relation.

For the probability-valid periodic mixture of
Theorem~\ref{thm:r72-periodic-memory}, with distinct nontrivial frequency
orders $q_1,\ldots,q_r$, let
\[
 K=\min_{\mathcal P}\sum_{B\in\mathcal P}
                  \operatorname{lcm}\{q_j:j\in B\},
\]
where $\mathcal P$ ranges over partitions of the frequency indices.
The exact pure-start count is the minimum of the full period
$\operatorname{lcm}(q_1,\ldots,q_r)$ and $K+1$, with the constant and
single-frequency endpoints as stated in the retained theorem.
A cyclic permutation representation is a disjoint union of cycles; a
frequency of order $q$ occurs on a cycle precisely when $q$ divides its
length. Assigning required modes to cycles therefore gives the same
partition expression. The bounded decoders and positive-amplitude
preparation in the retained theorem attain that lower bound. This agreement
identifies a model-specific closed form inside the general factorization
problem. Arbitrary bounded decoders do not have this same closed formula,
as both the sharp $S_3$ and the five-character contrast transitions show.
The quotient theorem applies to arbitrary finite abelian tables; it does
not assume the positive-amplitude mixture of the periodic closed form.

The leakage theorem, Theorem~\ref{thm:r72-periodic-leakage}, adds its
separately proved inactive live label. Its lower bound uses a
decay-normalized observable core and the no-critical-chain argument;
the common-background plateau of Section~\ref{sec:r73-finite-certificates}
is a different perturbation, as is terminal readout contrast.
The exact four-period damping threshold was already $1/\sqrt3$ by
Theorem~\ref{thm:r69-sector-phase}; the current supplement makes this
inherited closure explicit. The exact spectral construction of
Theorem~\ref{thm:r69-spectral}, the relative-spectrum obstruction and
the arbitrary-phase envelope of Theorem~\ref{thm:r71-envelope} remain
proved with their own hypotheses. In particular a strictly damped
irrational mode need not have compactly reversible observable dynamics.

The following actual-history transfer theorem remains a general statement
about causal experiment comparison. It compares complete procedures,
charges the common paid acquisition construction, and forms its lower
witness on actual competing histories. None of its hypotheses is inferred
from a finite observable group. The complete paired supplement retains the
associated performed-law minimax and dual theory, finite-width transfer,
reset control, optimal discrimination, predictive quotients and geometric
results. The original ordered-measurement and two-stage-reset realizations
remain independently identifiable by source commit and path.

\section{The causal task-order transfer theorem}\label{sec:central}
Suppose that every terminating state-stopped competitor reaches at least $B_0\ge1$ live labels under $\theta_*$. Such a certificate must concern actual physical command compatibility, not printed action names. Define
\begin{equation}\label{eq:widths}
\begin{array}{c|cc}
 i&m_-(M)&m_+(M)\\\hline
 \cp&M&\lfloor M/D_0\rfloor\\
 \ev&M&\lfloor M/D_0\rfloor\\
 \st&M-B_0&\lfloor(M-B)/D_0\rfloor .
\end{array}
\end{equation}
The event upper also requires $M\ge B$. The trivial state-stopped lower $B_0=1$ is valid for a procedure initially waiting to make a call.

\begin{theorem}[Affine task-order transfer with a paid reference and common construction]\label{thm:transfer}
Assume all the prepared-kernel, full-report, task-order, effort, occupation, resource and common-construction hypotheses of Section~\ref{sec:data}. Fix $N,M$, put $b_-=N/\kappa_-$, $b_+=N/\kappa_+$, and suppose $b_-\ge1$ and $m_-(M)\ge1$. Set
\[
 t_N=(1-A_N)_+/u_-.
\]
If $t_N>0$, every one of the three risks satisfies
\begin{equation}\label{eq:central-lower}
 R_i^\Theta(N,M)\ge
 \bigl[l_-t_N\Q_{m_-(M)}(b_-/t_N)-E_N\bigr]_+.
\end{equation}
If $t_N=0$, the asserted lower is zero, with no evaluation of $\Q$ at infinity. An empty legal class still has risk $+\infty$.

Suppose also $m_+(M)\ge1$, $b_+\ge1$, and, for event stopping, $M\ge B$. If
\begin{equation}\label{eq:beta}
 \beta_N=\frac{l_+b_+}{1+b_+a_+}\ge1,
\end{equation}
one common Borel program with actual expected cost at most $N$ in every model attains
\begin{equation}\label{eq:central-upper}
 R_i^\Theta(N,M)\le
 \frac{u_+(1+b_+a_+)}{l_+}\Q_{m_+(M)}(\beta_N)+b_+e_+.
\end{equation}
It commits almost surely. If $\beta_N<1$, the always-accept first-attempt construction instead gives $R_i^\Theta(N,M)\le u_+\Q_{m_+(M)}(1)+e_+$ at cost at most $\kappa_+\le N$.

The lower is against every legal competing strategy and uses one actual witness. The upper is one parameter-independent strategy for all models. The two effort constants need not agree. Neither direction assumes that the comparison experiment belongs to $\Theta$ or that a normalized predictive update is uniformly contracting.
\end{theorem}

\begin{corollary}[Scale-uniform relative transfer]\label{cor:relative}
Under Theorem~\ref{thm:transfer}, suppose $A_N=E_N=a_+=e_+=0$. If $l_+b_+\ge1$, then
\begin{equation}\label{eq:relative}
 \frac{l_-}{u_-}\Q_{m_-(M)}(u_-b_-)
 \le R_i^\Theta(N,M)
 \le\frac{u_+}{l_+}\Q_{m_+(M)}(l_+b_+).
\end{equation}
The factors $l_\pm,u_\pm$ may be fixed nonunit constants, independent of budget and resolution. In particular no term proportional to $b_\pm$ times a fixed absolute calibration error is required in this relative-order class.
\end{corollary}
\begin{proof}
Substitute $t_N=1/u_-$ and $\beta_N=l_+b_+$ in the theorem.
\end{proof}

\begin{corollary}[Separate exact checkpoint and autonomous laws]\label{cor:exact}
Suppose the certificates are exact: $l_\pm=u_\pm=1$, all defects vanish, $\kappa_-=\kappa_+=\kappa$, and exports are the identity. Then, for $N\ge\kappa$,
\[
 R_\cp^\Theta(N,m)=\Q_m(N/\kappa).
\]
The constructed autonomous upper uses $B+m$ state-stopped labels, or $\max(B,m)$ event-stopped labels. If the mandatory live count is sharp, $B_0=B$, then
\begin{align}
 R_\st^\Theta(N,B+m)&=\Q_m(N/\kappa)\quad(m\ge1),\label{eq:typed-exact}\\
 R_\ev^\Theta(N,M)&=\Q_M(N/\kappa)\quad(M\ge B).\label{eq:event-exact}
\end{align}
The subtraction by $B$ belongs only to the state-stopped model. Nonidentity exports require the expanded dictionary in \eqref{eq:widths}, unless an optimal dictionary is invariant under all exports. These exact stopping distinctions are inherited from the complete companion and are reproved here to fix the resource semantics.
\end{corollary}
\begin{proof}
The same lower and upper profile coincide in each asserted interface. Identity exports give $D_0=1$. The packing and converse widths are proved in Lemma~\ref{lem:roles}. No lower on live states is used to forbid event-stopped reuse.
\end{proof}

The familiar absolute-defect theorem is also a specialization: take $l_\pm=u_\pm=1$ while retaining $A_N,E_N,a_+,e_+$. Thus relative and absolute calibration are composed in one statement. The chosen metric is the common scored task throughout; it is not changed between the converse and the construction.

\section{Proof of the central theorem}
\begin{lemma}[Profile attainment]\label{lem:attainment}
For every finite $m\ge1$ and $b\ge1$, the minimum in \eqref{eq:profile} is attained. The profile is nonincreasing in $m,b$ and left-continuous in $b>1$.
\end{lemma}
\begin{proof}
Write a dictionary as an ordered tuple in the compact set $H^m$, permitting repetitions. The feasible densities are a weak-star closed subset of the radius-$b$ ball of $L^\infty(\nu)$: positivity, the upper bound and the mass equality are each closed under pairing with $L^1(\nu)$. They are therefore weak-star compact. As the tuple varies, its squared-distance function varies uniformly on $H$, since all points belong to a bounded set. Along a convergent product net, split the objective difference into this uniform variation and the weak-star pairing against the fixed bounded measurable loss at the limit tuple. This proves continuity of the objective and attainment. Nets avoid an unwarranted sequential compactness assumption. Monotonicity follows from inclusion of feasible sets. For left continuity take an optimizer $h$ at $b>1$ and let $h_\lambda=(1-\lambda)h+\lambda$. It has unit mass and cap $b_\lambda=(1-\lambda)b+\lambda<b$. Its objective tends to the optimum as $\lambda\downarrow0$, while monotonicity bounds the optimum at $b_\lambda$ between the optimum at $b$ and this objective.
\end{proof}

\begin{lemma}[Actual-history comparison measure]\label{lem:comparison}
Fix any robustly admissible competitor of cost at most $N$, with its fixed terminal dictionary $C$. There are finite nonnegative measures $\xi_c$ on $X$, $c\in C$, such that $\xi=\sum_c\xi_c$ obeys
\begin{align}
 \xi&\le (\E_*A_\tau)\nu\le b_-\nu,\label{eq:measure-cap}\\
 1&\le u_-\xi(X)+A_N,\label{eq:measure-mass}\\
 \E_*\|Z-\widehat r\|^2
 &\ge l_-\sum_{c\in C}\int d_c\,d\xi_c-E_N.
 \label{eq:measure-risk}
\end{align}
The measure $\xi$ is not asserted to be the terminal law of any comparison execution.
\end{lemma}
\begin{proof}
Adjoin an independent uniform random variable to the current report on both sides. This is current-call randomness, not a seed retained between calls; it is sufficient to represent the competitor's randomized output and stopping decision. A finite-output kernel has such a representation by cumulative probabilities. At a fixed reached productive history $h$ and chosen command $a$, the post-report decision then gives an acceptance indicator $q_h(y)$ and an output $r_h(y)\in C$ on acceptance. Give $r_h$ an arbitrary fixed value on rejection. The notation $y$ now includes the current independent random variable. Integrating the original certificates over that variable preserves their constants. Dependence on the entire earlier history is permitted in these functions.

For each Borel $D\subset X$, set
\begin{equation}\label{eq:xi-definition}
 \xi_c(D)=\E_*\sum_{j\le\tau:\ j\ {
m productive}}
 \int\one_{\{\pi_{H_{j-1},a_j}(y)\in D\}}
 q_{H_{j-1}}(y)\one_{\{r_{H_{j-1}}(y)=c\}}
 P^0_{H_{j-1},a_j}(dy,H).
\end{equation}
The outer expectation is the actual pre-call occupation at $\theta_*$. All terms are jointly measurable and nonnegative. Tonelli defines the measures, whose total is finite because $\E_*A_\tau\le b_-$. Summing over $c$ and dropping the acceptance indicator gives \eqref{eq:measure-cap}: each comparison row has mark law $\nu$.

The analogous sum with the actual rows instead of $P^0$ has mass exactly one, because the actual procedure commits exactly once almost surely. Apply \eqref{eq:lower-mass} at each actual reached history, then sum. The occupation certificate yields \eqref{eq:measure-mass}. Apply \eqref{eq:lower-task} in the same way to the disjoint actual commitment events. Their expected losses sum to the actual risk. The comparison contribution is $d_{r_h(y)}(\pi_{h,a}y)$ by \eqref{eq:full-moments}; this identity conditions on the full current report and therefore does not overlook extra information seen by the competitor. The accumulated loss defect is at most $E_N$, giving \eqref{eq:measure-risk}. One may first truncate the number of calls. Monotone convergence for the nonnegative contributions and the finite defect bounds justify the limit.

No actual past is replaced by a comparison past. In particular the competitor need not terminate, or obey a cost bound, if hypothetically run under the comparison experiment. This distinction is why the argument applies when that experiment is not an operating model.
\end{proof}

\begin{remark}[Why the full report is required]
A constant mark can conceal a visible fair bit which also equals the future audit. A comparison using only that mark has $f=0,w=1$, but a two-value decoder reading the bit predicts it with zero loss. The hypotheses therefore condition on the full report seen by the competitor; matching a compressed mark marginal alone cannot establish the lower.
\end{remark}

\begin{lemma}[Live roles, output roles and physical compatibility]\label{lem:roles}
A program with $\ell$ live roles and $t$ output roles has a state-stopped implementation with $\ell+t$ labels and an event-stopped implementation with $\max(\ell,t)$ labels, preserving every path cost and complete stopped scored law. Conversely an event-stopped $M$-label program has at most $M$ decoder values. A state-stopped $M$-label program with at least $B_0$ reached live labels has at most $M-B_0$ terminal decoder values.

For a physical live-state certificate, consider finitely many mandatory pre-call contexts $I$ with no free context input. Suppose the commands legal at a history of context $i$ are exactly $\mathcal U_i$, and every terminating competitor visits every context with positive probability under the same witness. This positive probability may depend on the competitor; no uniform lower bound over competitors is required. Let $b_{\rm phys}$ be the least size of a cover of $I$ by subsets $J$ with $\bigcap_{i\in J}\mathcal U_i\ne\varnothing$. Then at least $b_{\rm phys}$ live labels are required. This certificate is invariant under duplicate names and bijective stateless recodings of the same physical commands.
\end{lemma}
\begin{proof}
For state stopping take disjoint copies of the live and terminal role sets and type them zero and one. For event stopping embed the live and terminal role sets separately into one alphabet of size $\max(\ell,t)$. On continuing transitions use the live embedding; on commitment emit the final event and use the terminal embedding. A command row is read only on a continuing branch, while the decoder is read only after commitment. Sharing their label therefore causes no conflict and uses no continuing mode bit. Couple the original and recoded coins. Every finite stopped prefix and every path cost agree, and almost-sure termination gives equality of the stopped scored laws. A fixed decoder has one value per retained label. In state stopping the global halt map excludes each reached live label from the output set, proving its subtraction. Event stopping has no such exclusion.

A live label reached in contexts $J$ must have one command distribution supported on all the corresponding $\mathcal U_i$. Since $J$ is finite, the union of the exceptional null command sets is null; hence their intersection supports a probability and is nonempty. The sets of contexts served by the labels cover $I$, proving the bound. Conversely a nonempty intersection contains a command and hence a Dirac command distribution. Replacing names without changing executed kernels, legal domains, costs and reports does not change these intersections. A translator which itself remembers a hidden phase is not such a recoding; its memory must be charged.
\end{proof}

\begin{proof}[Proof of Theorem~\ref{thm:transfer}]
Fix a competing program. Lemma~\ref{lem:roles} bounds its decoder width by $m_-(M)$. Let $s=\xi(X)$ in Lemma~\ref{lem:comparison}. If $t_N>0$, then \eqref{eq:measure-mass} gives $s\ge t_N$. Scale the measure to $\xi'=(t_N/s)\xi$. It has mass $t_N$ and remains bounded above by $b_-\nu$. Thus $h=(d\xi'/d\nu)/t_N$ has unit mass and cap $b_-/t_N$. Nearest-dictionary projection and nonnegativity imply
\[
 \sum_{c\in C}\int d_c\,d\xi_c
 \ge\int [w+\dist(f,C)^2]\,d\xi
 \ge t_N\Q_{m_-(M)}(b_-/t_N).
\]
Together with \eqref{eq:measure-risk} this proves the asserted lower for this competitor. Its worst-model risk dominates its risk at $\theta_*$. Taking the infimum over competitors gives \eqref{eq:central-lower}, without exchanging an infimum and a supremum. Nonnegative loss gives the zero truncation and the stated case $t_N=0$.

For the upper assume \eqref{eq:beta}. Choose an attained optimizer $(C,h)$ at width $m_+(M)$ and cap $\beta_N$. Let $r_C(x)$ be the first nearest point to $f(x)$ in a fixed ordering of $C$. After each attempt of $\mathcal G$, accept $x$ with probability $h(x)/\beta_N$, export $T_Jr_C(x)$ and commit; otherwise return to the start routing label. All these operations are Borel. The choice $J$ is used within the current attempt; any part retained thereafter is charged through the live role or exported decoder, not a free seed. The exports are assumed jointly measurable and take at most $D_0|C|$ distinct values, so their finite decoder and transitions are Borel as well.

Conditional on every surviving attempt history, the acceptance probability is at least
\begin{equation}\label{eq:hazard}
 p=\frac{l_+}{\beta_N}-a_+=\frac1{b_+}>0.
\end{equation}
Consequently the probability of reaching attempt $j+1$ is at most $(1-p)^j$, by induction on conditional survival probabilities. In particular commitment is almost sure and the expected number of attempts is at most $1/p=b_+$. Each begun attempt has conditional expected cost at most $\kappa_+$, not merely a cost conditional on a favorable outcome. Summing the predictable attempt costs gives actual expected cost at most $\kappa_+b_+=N$.

The conditional unnormalized loss contributed by an attempt is at most $u_+\Q_{m_+(M)}(\beta_N)/\beta_N+e_+$. Sum its contributions on the disjoint terminal events, using the same survival bound. The result is at most
\[
 b_+\left[\frac{u_+}{\beta_N}\Q_{m_+(M)}(\beta_N)+e_+\right],
\]
which is \eqref{eq:central-upper}. The program has $B$ live roles and at most $D_0m_+(M)$ terminal roles. Lemma~\ref{lem:roles} gives exactly the sufficient widths in \eqref{eq:widths}. A checkpoint counts only the terminal width. No physical renewal law was used: the conditional hazard controls survival even if the hidden state after rejection differs arbitrarily from its former state.

If $\beta_N<1$, commit after the first attempt, and use an optimizer at cap one. The loss certificate with $q=1$ gives $u_+\Q_{m_+(M)}(1)+e_+$, and the whole-attempt cost is at most $\kappa_+$. This proves all remaining assertions.
\end{proof}

\section{Scope and further mathematical questions}\label{sec:r73-scope}\label{sec:r74-scope}
The compactly reversible normal form is a classification of exact
input--response realization. The contrast theorem applies to every finite
observable group, while the quotient linear programs apply to groups whose
subgroups are normal. The sharp $S_3$ and $(\mathbb F_2)^4$ thresholds are
complete solutions of specified models. They do not identify every
nonabelian positive-factorization minimum. The full contrast filtration is
finite and algebraically computable for algebraic data; a closed formula
for all groups or a polynomial synthesis algorithm remains a different
question.

The retained posterior-positivity theorem handles arbitrary finite collections of
known instruments satisfying its uniform hypothesis, including
state-dependent observations and feedback from visible histories. It gives
one exact simulator for a fixed known initial state. It does not convert an
unknown quantum input into a classical label, learn unknown primitives, or
produce parameter-independent tables for an arbitrary statistical model.
The retained geometric lower bound is for a physical dictionary
covering a whole interior state body. Its matching small-noise exponent is
not a lower bound for the unrestricted memory of each individual instrument.
The rank-three family instead has exact unrestricted memory lower bounds,
including the separate charge for pure initialization.

Terminal readout contrast, channel damping repeated after each command,
and noise on every observed instrument branch are different operations.
The four-period damping transition at $1/\sqrt3$ was already determined by
the retained first-sector theorem; Corollary~\ref{cor:r74-period-four-closed}
records that consequence explicitly. The previous finite-probe gap remains
valid. Other damping transitions, intermediate phase-type timing frontiers
and the general $M$ squared-sign risk curve are not settled here. The exact
periodic partition and leakage laws keep their spectral and probability
hypotheses. The arbitrary-phase arithmetic envelope keeps both logarithmic
exponents and its badly-approximable criterion; it supplies no particular
irrationality exponent for the algebraic phase $(3+4i)/5$.

Full-history optimization, finite-memory attainment, finite-memory density,
approximation at a fixed memory budget and exact response simulation are
distinct problems. The robust tied criterion already proved on the rigid
herald class remains in force. General nonrigid boundary ties and a
classification of optimal several-command feedback remain open. Exact
algebraic descriptions of transition tables also leave sampling, precision,
program length and working storage as additional resources.

The actual-history transfer theorem and the performed-law, variational,
dual, spectral, reset, discrimination and geometric results retain complete
proofs. The independent absolute-calibration companion and the original
finite-order measurement and two-stage reset realizations retain their
source identities. Claims requiring A1, A2 or further Theta construction
hypotheses remain independent research objectives. The scientific question
throughout is how the information acquired by a causal experiment constrains
the resources required to realize and use its predictive state.
\begin{thebibliography}{99}
\bibitem{Blackwell} D. Blackwell, \emph{Equivalent comparisons of experiments}, Ann. Math. Statist. \textbf{24} (1953), 265--272. DOI: 10.1214/aoms/1177729032.
\bibitem{CGM97} J. A. Cuesta-Albertos, A. Gordaliza and C. Matr\'an, \emph{Trimmed $k$-means: an attempt to robustify quantizers}, Ann. Statist. \textbf{25} (1997), 553--576. DOI: 10.1214/aos/1031833664.
\bibitem{GG99} L. A. Garc\'ia-Escudero and A. Gordaliza, \emph{Robustness properties of $k$ means and trimmed $k$ means}, J. Amer. Statist. Assoc. \textbf{94} (1999), 956--969. DOI: 10.1080/01621459.1999.10474200.
\bibitem{SC} C. R. Shalizi and J. P. Crutchfield, \emph{Computational mechanics: pattern and prediction, structure and simplicity}, J. Statist. Phys. \textbf{104} (2001), 817--879. DOI: 10.1023/A:1010388907793.
\bibitem{BC} N. Barnett and J. P. Crutchfield, \emph{Computational mechanics of input--output processes: structured transformations and the $\varepsilon$-transducer}, J. Statist. Phys. \textbf{161} (2015), 404--451; arXiv:1412.2690.
\bibitem{AMTW} A. Ambainis, M. Mosca, A. Tapp and R. de Wolf, \emph{Private quantum channels}, Proc. 41st Annual Symposium on Foundations of Computer Science (2000), 547--553. DOI: 10.1109/SFCS.2000.892142.
\bibitem{Buscemi} F. Buscemi, \emph{Comparison of quantum statistical models: equivalent conditions for sufficiency}, Comm. Math. Phys. \textbf{310} (2012), 625--647. DOI: 10.1007/s00220-012-1421-3.
\bibitem{BR14} N. Bertschinger and J. Rauh, \emph{The Blackwell relation defines no lattice}, Proc. IEEE International Symposium on Information Theory (2014), 2479--2483. DOI: 10.1109/ISIT.2014.6875280; arXiv:1401.3146.
\bibitem{Banerjee05} A. Banerjee, X. Guo and H. Wang, \emph{On the optimality of conditional expectation as a Bregman predictor}, IEEE Trans. Inform. Theory \textbf{51} (2005), 2664--2669. DOI: 10.1109/TIT.2005.850145.
\bibitem{Stengel96} B. von Stengel, \emph{Efficient computation of behavior strategies}, Games Econom. Behav. \textbf{14} (1996), 220--246. DOI: 10.1006/game.1996.0050.
\bibitem{Balder89} E. J. Balder, \emph{On compactness of the space of policies in stochastic dynamic programming}, Stochastic Process. Appl. \textbf{32} (1989), 141--150. DOI: 10.1016/0304-4149(89)90058-6.
\bibitem{Barry14} J. Barry, D. T. Barry and S. Aaronson, \emph{Quantum partially observable Markov decision processes}, Phys. Rev. A \textbf{90} (2014), 032311. DOI: 10.1103/PhysRevA.90.032311; arXiv:1406.2858.
\bibitem{Ying18} S. Ying and M. Ying, \emph{Reachability analysis of quantum Markov decision processes}, Inform. Comput. \textbf{263} (2018), 31--51. DOI: 10.1016/j.ic.2018.09.001; arXiv:1406.6146.
\bibitem{Saldi25} N. Saldi, S. Sanjari and S. Y\"uksel, \emph{Quantum Markov decision processes: general theory, approximations, and classes of policies}, SIAM J. Control Optim. \textbf{63} (2025), 1156--1183. DOI: 10.1137/24M1643864; arXiv:2402.14649.
\bibitem{Seidenberg54} A. Seidenberg, \emph{A new decision method for elementary algebra}, Ann. of Math. (2) \textbf{60} (1954), 365--374. DOI: 10.2307/1969640.
\bibitem{Chatterjee15} K. Chatterjee, M. Chmel\'ik, R. Gupta and A. Kanodia, \emph{Optimal cost almost-sure reachability in POMDPs}, Proc. AAAI Conference on Artificial Intelligence \textbf{29} (2015). DOI: 10.1609/aaai.v29i1.9683.
\bibitem{Andriushchenko23} R. Andriushchenko, A. Bork, M. \v{C}e\v{s}ka, S. Junges, J.-P. Katoen and M. Mac\'ak, \emph{Search and explore: symbiotic policy synthesis in POMDPs}, Computer Aided Verification (2023); arXiv:2305.14149v2.
\bibitem{Halmos69} P. R. Halmos, \emph{Two subspaces}, Trans. Amer. Math. Soc. \textbf{144} (1969), 381--389. DOI: 10.1090/S0002-9947-1969-0251519-5.
\bibitem{Kessler21} D. A. Kessler, E. Barkai and K. Ziegler, \emph{First-detection time of a quantum state under random probing}, Phys. Rev. A \textbf{103} (2021), 022222. DOI: 10.1103/PhysRevA.103.022222; arXiv:2012.01763.
\bibitem{DelVecchio26} G. Del Vecchio Del Vecchio and S. N. Majumdar, \emph{Optimal detection of quantum states via projective measurements}, J. Phys. A: Math. Theor. \textbf{59} (2026), 035001. DOI: 10.1088/1751-8121/ae34fe; arXiv:2509.08556.
\bibitem{Companion} Q. Qi, \emph{General Theta Foundations I: Acquired Geometry and Causal Resource Transfer}, complete technical companion on absolute calibration, exact stopping resources and nonintervening sensing, 2026. The complete retained R58 native source and its independently built companion accompany this manuscript; their identities are fixed in the source-provenance and preservation records.
\bibitem{Tzeng92} W.-G. Tzeng, \emph{A polynomial-time algorithm for the equivalence of probabilistic automata}, SIAM J. Comput. \textbf{21} (1992), 216--227. DOI: 10.1137/0221017.
\bibitem{LiYuYing12} Y. Li, N. Yu and M. Ying, \emph{Termination of nondeterministic quantum programs}, arXiv:1201.0891 (2012).
\bibitem{Fu24} J. Fu, H. Jiang, M. Xu, Y. Deng and Z.-B. Li, \emph{Algorithmic analysis of termination problems for nondeterministic quantum programs}, arXiv:2402.15827 (2024).

\bibitem{LiuVerdu18} J. Liu and S. Verd\'u, \emph{Rejection sampling and noncausal sampling under moment constraints}, Proc. IEEE International Symposium on Information Theory (2018), 1565--1569. DOI: 10.1109/ISIT.2018.8437857.
\bibitem{GocFlamich24} D. Goc and G. Flamich, \emph{On channel simulation with causal rejection samplers}, Proc. IEEE International Symposium on Information Theory (2024). DOI: 10.1109/ISIT57864.2024.10619339; examined text arXiv:2401.16579v2 (3 May 2024).
\bibitem{LiLiu26} Z. Li and T. Liu, \emph{Ranking experiments under sequential sampling}, arXiv:2608.19897v2 (28 August 2026). Version 1 appeared under the title \emph{The order of binary experiments under endogenous stopping}, with Z. Li as sole author; the version-2 statement is the one compared here.
\bibitem{Ying14pre} S. Ying and M. Ying, \emph{Reachability analysis of quantum Markov decision processes}, arXiv:1406.6146 (2014), Theorem 3.6; published journal version listed above as \cite{Ying18}.

\bibitem{KaraYuksel20} A. D. Kara and S. Y\"uksel, \emph{Robustness to incorrect system models in stochastic control}, SIAM J. Control Optim. \textbf{58} (2020), 1144--1182. DOI: 10.1137/18M1208058; examined text arXiv:1803.06046v7.
\bibitem{Demirci25} Y. E. Demirci, A. D. Kara and S. Y\"uksel, \emph{Sensitivity of filter kernels and robustness bounds to transition and measurement kernel perturbations in partially observable stochastic control}, arXiv:2508.10658v3 (2025).
\bibitem{Watrous18} J. Watrous, \emph{The Theory of Quantum Information}, Cambridge University Press, 2018, Chapter 3, Corollary 3.40.
\bibitem{KRN65} K. Kuratowski and C. Ryll-Nardzewski, \emph{A general theorem on selectors}, Bull. Acad. Polon. Sci. S\'er. Sci. Math. Astronom. Phys. \textbf{13} (1965), 397--403.

\bibitem{Charnes63} A. Charnes, W. W. Cooper and K. O. Kortanek, \emph{Duality in semi-infinite programs and some works of Haar and Carath\'eodory}, Management Sci. \textbf{9} (1963), 209--228. DOI: 10.1287/mnsc.9.2.209.
\bibitem{Shapiro09} A. Shapiro, \emph{Semi-infinite programming, duality, discretization and optimality conditions}, Optimization \textbf{58} (2009), 133--161. DOI: 10.1080/02331930902730070; author preprint, Optimization Online (2008).
\bibitem{Dufour25} F. Dufour and T. Prieto-Rumeau, \emph{Absorbing Markov decision processes and their occupation measures}, SIAM J. Control Optim. \textbf{63} (2025), 676--698. DOI: 10.1137/24M1661959. Related author preprint examined: \emph{A note on weak compactness of occupation measures for an absorbing Markov decision process}, arXiv:2402.10672v2 (2024).
\bibitem{Franc23} V. Franc, D. Prusa and V. Voracek, \emph{Optimal strategies for reject option classifiers}, J. Mach. Learn. Res. \textbf{24} (2023), no. 11, 1--49.
\bibitem{Jencova16} A. Jen\v{c}ov\'a, \emph{Comparison of quantum channels and statistical experiments}, arXiv:1512.07016v2 (2016). The separately indexed shorter version arXiv:1601.06370 was withdrawn in favor of this extended text.
\bibitem{LiElGamal18} C. T. Li and A. El Gamal, \emph{Strong functional representation lemma and applications to coding theorems}, IEEE Trans. Inform. Theory \textbf{64} (2018), 6967--6978; arXiv:1701.02827v4.
\bibitem{AldousShepp87} D. Aldous and L. Shepp, \emph{The least variable phase type distribution is Erlang}, Comm. Statist. Stochastic Models \textbf{3} (1987), 467--473. DOI: 10.1080/15326348708807067.
\bibitem{Sion58} M. Sion, \emph{On general minimax theorems}, Pacific J. Math. \textbf{8} (1958), 171--176. DOI: 10.2140/pjm.1958.8.171.
\bibitem{PiunovskiyZhang24} A. Piunovskiy and Y. Zhang, \emph{Extreme occupation measures in Markov decision processes with an absorbing state}, SIAM J. Control Optim. \textbf{62} (2024), 65--90. DOI: 10.1137/23M1572398. Author text examined: \emph{Extreme occupation measures in Markov decision processes with a cemetery}, arXiv:2307.03158v1.
\bibitem{DufourMixtures} F. Dufour and T. Prieto-Rumeau, \emph{Absorbing Markov decision processes: geometric properties and sufficiency of finite mixtures of deterministic policies}, arXiv:2512.17511; examined HTML at version-1 URL, Theorems 4.3--4.4 and Proposition 4.5.
\bibitem{Mereghetti24} C. Mereghetti, B. Palano and P. Raucci, \emph{Unary quantum finite state automata with control language}, Appl. Sci. \textbf{14} (2024), no. 4, 1490. DOI: 10.3390/app14041490.
\bibitem{RabinScott59} M. O. Rabin and D. Scott, \emph{Finite automata and their decision problems}, IBM J. Res. Develop. \textbf{3} (1959), 114--125. DOI: 10.1147/rd.32.0114.
\bibitem{Jaeger00} H. Jaeger, \emph{Observable operator models for discrete stochastic time series}, Neural Comput. \textbf{12} (2000), 1371--1398. DOI: 10.1162/089976600300015411.
\bibitem{ThonJaeger15} M. Thon and H. Jaeger, \emph{Links between multiplicity automata, observable operator models and predictive state representations---a unified learning framework}, J. Mach. Learn. Res. \textbf{16} (2015), 103--147.
\bibitem{MonrasWinter14} A. Monras and A. Winter, \emph{Quantum learning of classical stochastic processes: the completely-positive realization problem}, LIPIcs \textbf{27}, TQC 2014 (2014), 99--109. DOI: 10.4230/LIPIcs.TQC.2014.99.

\bibitem{Heller65} A. Heller, \emph{On stochastic processes derived from Markov chains}, Ann. Math. Statist. \textbf{36} (1965), 1286--1291. DOI: 10.1214/aoms/1177700000.
\bibitem{CarlylePaz71} J. W. Carlyle and A. Paz, \emph{Realizations by stochastic finite automata}, J. Comput. System Sci. \textbf{5} (1971), 26--40. DOI: 10.1016/S0022-0000(71)80005-3.
\bibitem{Fanizza24} M. Fanizza, J. Lumbreras and A. Winter, \emph{Quantum theory in finite dimension cannot explain every general process with finite memory}, Comm. Math. Phys. \textbf{405} (2024), article 50. DOI: 10.1007/s00220-023-04913-4; arXiv:2209.11225. See Section 2.3 for the FRDN phase examples and Theorem 2 and Section 4.2 for the noise-dependent worst-phase classical-memory lower bound.
\bibitem{DmitrievDynkin46} N. Dmitriev and E. Dynkin, \emph{On characteristic roots of stochastic matrices}, Izv. Akad. Nauk SSSR Ser. Mat. \textbf{10} (1946), no. 2, 167--184. Math-Net.Ru: im3595 (Russian).
\bibitem{JohnsonPaparella17} C. R. Johnson and P. Paparella, \emph{A matricial view of the Karpelevi\v c theorem}, Linear Algebra Appl. \textbf{520} (2017), 1--15. DOI: 10.1016/j.laa.2017.01.009; arXiv:1611.06970v2.
\bibitem{KirklandSmigoc22} S. Kirkland and H. \v Smigoc, \emph{Stochastic matrices realising the boundary of the Karpelevi\v c region}, Linear Algebra Appl. \textbf{635} (2022), 116--138. DOI: 10.1016/j.laa.2021.11.016; arXiv:2110.01040v1.
\bibitem{BenvenutiFarina04} L. Benvenuti and L. Farina, \emph{A tutorial on the positive realization problem}, IEEE Trans. Automat. Control \textbf{49} (2004), 651--664. DOI: 10.1109/TAC.2004.826715.
\bibitem{BenvenutiFarina99} L. Benvenuti and L. Farina, \emph{An example of how positivity may force realizations of ``large'' dimension}, Systems Control Lett. \textbf{36} (1999), 261--266. DOI: 10.1016/S0167-6911(98)00098-X.
\bibitem{OCinneide90} C. A. O'Cinneide, \emph{Characterization of phase-type distributions}, Comm. Statist. Stochastic Models \textbf{6} (1990), 1--57. DOI: 10.1080/15326349908807134.
\bibitem{Maier91} R. S. Maier, \emph{The algebraic construction of phase-type distributions}, Comm. Statist. Stochastic Models \textbf{7} (1991), 573--602. DOI: 10.1080/15326349108807207.
\bibitem{HoneVarona19} A. N. W. Hone and J. L. Varona, \emph{Continued fractions and irrationality exponents for modified Engel and Pierce series}, Monatsh. Math. \textbf{190} (2019), 501--516. DOI: 10.1007/s00605-018-1244-1. Author version arXiv:1810.10127v1, equation (3.3).
\bibitem{Nussbaum98} R. D. Nussbaum and M. Scheutzow, \emph{Admissible arrays and a nonlinear generalization of Perron--Frobenius theory}, J. London Math. Soc. (2) \textbf{58} (1998), 526--544. DOI: 10.1112/S0024610798006796.
\bibitem{Brown64} D.~R. Brown, On clans of non-negative matrices,
\emph{Proc. Amer. Math. Soc.} \textbf{15} (1964), 671--674,
doi:10.1090/S0002-9939-1964-0165504-6.
\bibitem{Schwarz64} \v{S}. Schwarz, On the structure of the semigroup of
stochastic matrices, \emph{Publ. Math. Inst. Hungar. Acad. Sci.}
\textbf{9} (1964), 297--311, especially Theorems~7--8.
\bibitem{Flor69} P. Flor, On groups of non-negative matrices,
\emph{Compositio Math.} \textbf{21} (1969), 376--382,
\url{https://www.numdam.org/article/CM_1969__21_4_376_0.pdf}.
\bibitem{Serre07} J.-P. Serre, Bounds for the orders of the finite subgroups
of $G(k)$, in \emph{Group Representation Theory}, M. Geck, D. Testerman
and J. Th\'evenaz (eds.), EPFL Press, Lausanne, 2007, 405--450;
author lecture text, April 2006, archived as arXiv:1011.0346.
\bibitem{BabaiBealsRockmore93} L. Babai, R. Beals and D. Rockmore,
Deciding finiteness of matrix groups in deterministic polynomial time,
\emph{Proceedings of ISSAC 1993}, 117--126,
doi:10.1145/164081.164104.
\bibitem{KieferWachter14} S. Kiefer and B. Wachter, Stability and complexity
of minimising probabilistic automata, in \emph{ICALP 2014, Part II},
Lecture Notes in Computer Science \textbf{8573}, 268--279,
doi:10.1007/978-3-662-43951-7\_23; arXiv:1404.6673, Section~4.2.
\bibitem{Schwarz62} \v{S}. Schwarz, Homomorphisms of a completely simple
semigroup onto a group, \emph{Matematicko-fyzik\'alny \v{c}asopis}
\textbf{12} (1962), no.~4, 293--300;
\url{https://www.dml.cz/handle/10338.dmlcz/126333}.

\bibitem{GPT13} J. Gouveia, P. A. Parrilo and R. R. Thomas,
\emph{Lifts of convex sets and cone factorizations}, Math. Oper. Res.
\textbf{38} (2013), 248--264. DOI: 10.1287/moor.1120.0575;
arXiv:1111.3164, especially Theorem 2.12 on equivariant factorizations.
\bibitem{KPT10} V. Kaibel, K. Pashkovich and D. O. Theis,
\emph{Symmetry matters for sizes of extended formulations},
SIAM J. Discrete Math. \textbf{26} (2012), 1361--1382.
DOI: 10.1137/110839813; arXiv:0911.3712v3.
An extended abstract appeared in IPCO 2010, LNCS 6080, 135--148.
\bibitem{BocciCarliniRapallo11} C. Bocci, E. Carlini and F. Rapallo,
\emph{Perturbation of matrices and nonnegative rank with a view toward
statistical models}, SIAM J. Matrix Anal. Appl. \textbf{32} (2011),
1500--1512. DOI: 10.1137/110825455; arXiv:1010.3566.
\bibitem{GillisGlineur12} N. Gillis and F. Glineur,
\emph{On the geometric interpretation of the nonnegative rank},
Linear Algebra Appl. \textbf{437} (2012), 2685--2712.
DOI: 10.1016/j.laa.2012.06.038; arXiv:1009.0880.
\bibitem{Okay26} C. Okay, \emph{Polyhedral classical simulators for quantum
computation}, Res. Math. Sci. \textbf{13} (2026), article 55.
DOI: 10.1007/s40687-026-00626-x. In particular Theorem 16, Lemma 17
and Corollary 19 treat invariant-polytope updates and adaptive composition.
\bibitem{AubrunSzarek17} G. Aubrun and S. J. Szarek,
\emph{Dvoretzky's theorem and the complexity of entanglement detection},
Discrete Anal. (2017), paper no. 1, 20 pp.; arXiv:1510.00578,
especially Section 6 on polyhedral approximations of density matrices.
\bibitem{Fawzi21} H. Fawzi,
\emph{On polyhedral approximations of the positive semidefinite cone},
Math. Oper. Res. \textbf{46} (2021), 1479--1489.
DOI: 10.1287/moor.2020.1077; arXiv:1811.09649.
\bibitem{HolroydLiggett16} A. E. Holroyd and T. M. Liggett,
\emph{Finitely dependent coloring}, Forum Math. Pi \textbf{4} (2016),
e9, 43 pp. DOI: 10.1017/fmp.2016.7; arXiv:1403.2448.
See Section 5 and the finite hidden-Markov obstruction attributed there
to Schramm.
\end{thebibliography}

\end{document}
